%%%%%%%%%%%%%%%%%%%%%%%%%%%%%%%%%%%%%%%%%%%%%%%%%%%%%%%%%%%%%%%%%%%%%%%
%% CATFISH-WY manuscript -- Oxford University Press authoring template
%% Build:  pdflatex catfish_wy -> bibtex catfish_wy -> pdflatex x2
%% Numbers in the Results come from tex/auto/numbers.tex, which is written
%% by scripts/07_tables_and_numbers.R after scripts 01-04 have been run.
%% Until then the red [run 07] placeholders mark every unfilled value.
%%%%%%%%%%%%%%%%%%%%%%%%%%%%%%%%%%%%%%%%%%%%%%%%%%%%%%%%%%%%%%%%%%%%%%%
\documentclass[unnumsec,webpdf,contemporary,large]{oup-authoring-template}

\theoremstyle{thmstyleone}
\newtheorem{theorem}{Theorem}
\newtheorem{proposition}[theorem]{Proposition}
\theoremstyle{thmstyletwo}
\newtheorem{example}{Example}
\newtheorem{remark}{Remark}
\theoremstyle{thmstylethree}
\newtheorem{definition}{Definition}

% give LaTeX room to place the full-width floats
\renewcommand{\topfraction}{0.92}
\renewcommand{\bottomfraction}{0.7}
\renewcommand{\textfraction}{0.08}
\renewcommand{\floatpagefraction}{0.75}
\renewcommand{\dbltopfraction}{0.92}
\renewcommand{\dblfloatpagefraction}{0.75}
\setcounter{topnumber}{3}
\setcounter{totalnumber}{5}
\setcounter{dbltopnumber}{3}

\usepackage{booktabs}
\usepackage{makecell}
\usepackage{multirow}
\usepackage{array}
\usepackage{tabularx}

\graphicspath{{../figures/}}

%% ---- placeholder / auto-number machinery ------------------------------
%% Placeholder definitions. Overwritten at \input time by auto/numbers.tex
%% once scripts 01-04 and 07 have been run.
\providecommand{\PHNA}{\textcolor{red}{[run~07]}}
\providecommand{\nTAnInd}{\PHNA}
\providecommand{\nTAnStr}{\PHNA}
\providecommand{\nTMvnInd}{\PHNA}
\providecommand{\nTMvnStr}{\PHNA}
\providecommand{\nTAnMax}{\PHNA}
\providecommand{\nTAnMin}{\PHNA}
\providecommand{\nTMvnMax}{\PHNA}
\providecommand{\nTMvnMin}{\PHNA}
\providecommand{\nTWyMax}{\PHNA}
\providecommand{\nTWyMin}{\PHNA}
\providecommand{\nLamAnMax}{\PHNA}
\providecommand{\nLamAnMin}{\PHNA}
\providecommand{\nLamMvnMax}{\PHNA}
\providecommand{\nLamMvnMin}{\PHNA}
\providecommand{\nTAnMaxOne}{\PHNA}
\providecommand{\nTMvnMaxOne}{\PHNA}
\providecommand{\nBrkAnKzero}{\PHNA}
\providecommand{\nBrkAnKfour}{\PHNA}
\providecommand{\nBrkMvnKzero}{\PHNA}
\providecommand{\nBrkMvnKfour}{\PHNA}
\providecommand{\nBrkMvnMax}{\PHNA}
\providecommand{\nBrkAnMax}{\PHNA}
\providecommand{\nLooFullMvn}{\PHNA}
\providecommand{\nLooMvnMin}{\PHNA}
\providecommand{\nLooMvnMax}{\PHNA}
\providecommand{\nLooWorst}{\PHNA}
\providecommand{\nLooWorstVal}{\PHNA}
\providecommand{\nRankAn}{\PHNA}
\providecommand{\nRankMvn}{\PHNA}
\providecommand{\nRandExp}{\PHNA}
\providecommand{\nNPath}{\PHNA}
\providecommand{\nNTrue}{\PHNA}
\providecommand{\nMu}{\PHNA}
\providecommand{\nWyThrMin}{\PHNA}
\providecommand{\nWyThrMax}{\PHNA}
\providecommand{\nBonfThr}{\PHNA}
\providecommand{\nFwerRatioMin}{\PHNA}
\providecommand{\nFwerRatioMax}{\PHNA}
\providecommand{\nBrkConsAnKfour}{\PHNA}
\providecommand{\nBrkConsMvnKfour}{\PHNA}
\providecommand{\nBrkMvnRangeLo}{\PHNA}
\providecommand{\nBrkMvnRangeHi}{\PHNA}
\providecommand{\nBrkMembership}{\PHNA}
\providecommand{\nBrkWyKzero}{\PHNA}
\providecommand{\nBrkBonfKzero}{\PHNA}
\providecommand{\nBrkWyKfour}{\PHNA}
\providecommand{\nBrkBonfKfour}{\PHNA}
\providecommand{\nBrkConsWyKfour}{\PHNA}
\providecommand{\nBrkWyRangeLo}{\PHNA}
\providecommand{\nBrkWyRangeHi}{\PHNA}
\providecommand{\nLooFullWy}{\PHNA}
\providecommand{\nLooWyMin}{\PHNA}
\providecommand{\nLooWyMax}{\PHNA}
\providecommand{\nLooThrRatio}{\PHNA}
\providecommand{\nLooBonfInd}{\PHNA}
\providecommand{\nLooBonfMod}{\PHNA}
\providecommand{\nLooBonfStr}{\PHNA}
\providecommand{\nLooBonfMax}{\PHNA}
\providecommand{\nLooWyInd}{\PHNA}
\providecommand{\nLooWyMod}{\PHNA}
\providecommand{\nLooWyStr}{\PHNA}
\providecommand{\nMissMvnZero}{\PHNA}
\providecommand{\nMissWyZero}{\PHNA}
\providecommand{\nMissMvnHalf}{\PHNA}
\providecommand{\nMissWyHalf}{\PHNA}
\providecommand{\nMissMvnFull}{\PHNA}
\providecommand{\nMissWyFull}{\PHNA}
\providecommand{\nMissMvnStrFull}{\PHNA}
\providecommand{\nMissWyStrFull}{\PHNA}
\providecommand{\nMissBonfMod}{\PHNA}
\providecommand{\nMissBonfStr}{\PHNA}
\providecommand{\nMissRatioMvnHalf}{\PHNA}
\providecommand{\nMissRatioWyHalf}{\PHNA}
\providecommand{\nMissLamLo}{\PHNA}
\providecommand{\nMissLamHi}{\PHNA}
\providecommand{\nMissCrossover}{\PHNA}
\providecommand{\nCmpWorst}{\PHNA}
\providecommand{\nCmpBest}{\PHNA}
\providecommand{\nCmpTFisherNA}{\PHNA}
\providecommand{\nCmpACATInd}{\PHNA}
\providecommand{\nCmpACATStr}{\PHNA}
\providecommand{\nCmpACATCal}{\PHNA}
\providecommand{\nCmpFisherInd}{\PHNA}
\providecommand{\nCmpFisherStr}{\PHNA}
\providecommand{\nCmpFisherCal}{\PHNA}
\providecommand{\nCmpStoufferInd}{\PHNA}
\providecommand{\nCmpStoufferStr}{\PHNA}
\providecommand{\nCmpStoufferCal}{\PHNA}
\providecommand{\nCmpminPInd}{\PHNA}
\providecommand{\nCmpminPStr}{\PHNA}
\providecommand{\nCmpminPCal}{\PHNA}
\providecommand{\nCmpTFisherInd}{\PHNA}
\providecommand{\nCmpTFisherStr}{\PHNA}
\providecommand{\nCmpTFisherCal}{\PHNA}
\providecommand{\nCmpOmnibusInd}{\PHNA}
\providecommand{\nCmpOmnibusStr}{\PHNA}
\providecommand{\nCmpOmnibusCal}{\PHNA}
\providecommand{\nCandN}{\PHNA}
\providecommand{\nCandNGwas}{\PHNA}
\providecommand{\nCandNPath}{\PHNA}
\providecommand{\nCandCtrlFound}{\PHNA}
\providecommand{\nCandCtrlTotal}{\PHNA}
\providecommand{\nCandCtrlMedian}{\PHNA}
\providecommand{\nCandCtrlTopPct}{\PHNA}
\providecommand{\nCandRandom}{\PHNA}
\providecommand{\nCandLDLR}{\PHNA}
\providecommand{\nCandPCSKnine}{\PHNA}
\providecommand{\nCandGwasTies}{\PHNA}

\IfFileExists{auto/numbers.tex}{%% AUTO-GENERATED by 07_tables_and_numbers.R -- do not edit by hand
\renewcommand{\nTAnInd}{0.0362}
\renewcommand{\nTAnStr}{0.3386}
\renewcommand{\nTMvnInd}{0.0472}
\renewcommand{\nTMvnStr}{0.0425}
\renewcommand{\nTAnMax}{0.3386}
\renewcommand{\nTAnMin}{0.0362}
\renewcommand{\nTMvnMax}{0.0568}
\renewcommand{\nTMvnMin}{0.0400}
\renewcommand{\nTWyMax}{0.0500}
\renewcommand{\nTWyMin}{0.0500}
\renewcommand{\nLamAnMax}{1.065}
\renewcommand{\nLamAnMin}{0.115}
\renewcommand{\nLamMvnMax}{1.181}
\renewcommand{\nLamMvnMin}{0.925}
\renewcommand{\nTAnMaxOne}{0.3011}
\renewcommand{\nTMvnMaxOne}{0.0157}
\renewcommand{\nBrkMembership}{8.8}
\renewcommand{\nBrkAnKzero}{0.0788}
\renewcommand{\nBrkMvnKzero}{0.0475}
\renewcommand{\nBrkWyKzero}{0.0547}
\renewcommand{\nBrkBonfKzero}{0.1290}
\renewcommand{\nBrkAnKfour}{0.1678}
\renewcommand{\nBrkMvnKfour}{0.0478}
\renewcommand{\nBrkWyKfour}{0.0549}
\renewcommand{\nBrkBonfKfour}{0.3111}
\renewcommand{\nBrkConsAnKfour}{0.0349}
\renewcommand{\nBrkConsMvnKfour}{0.0471}
\renewcommand{\nBrkConsWyKfour}{0.0529}
\renewcommand{\nBrkMvnMax}{0.0531}
\renewcommand{\nBrkAnMax}{0.2483}
\renewcommand{\nBrkWyRangeLo}{0.0420}
\renewcommand{\nBrkWyRangeHi}{0.0673}
\renewcommand{\nBrkMvnRangeLo}{0.0336}
\renewcommand{\nBrkMvnRangeHi}{0.0536}
\renewcommand{\nLooFullMvn}{0.0475}
\renewcommand{\nLooFullWy}{0.0547}
\renewcommand{\nLooMvnMin}{0.0326}
\renewcommand{\nLooMvnMax}{0.0539}
\renewcommand{\nLooWyMin}{0.0400}
\renewcommand{\nLooWyMax}{0.0680}
\renewcommand{\nLooWorst}{minP}
\renewcommand{\nLooWorstVal}{0.0543}
\renewcommand{\nLooThrRatio}{384.05}
\renewcommand{\nLooBonfInd}{0.0227}
\renewcommand{\nLooBonfMod}{0.1460}
\renewcommand{\nLooBonfStr}{0.1997}
\renewcommand{\nLooBonfMax}{0.2640}
\renewcommand{\nLooWyInd}{0.0587}
\renewcommand{\nLooWyMod}{0.0495}
\renewcommand{\nLooWyStr}{0.0537}
\renewcommand{\nRankAn}{11.8}
\renewcommand{\nRankMvn}{11.8}
\renewcommand{\nRandExp}{150}
\renewcommand{\nNPath}{300}
\renewcommand{\nNTrue}{15}
\renewcommand{\nMu}{1.5}
\renewcommand{\nWyThrMin}{\ensuremath{1.47\times10^{-4}}}
\renewcommand{\nWyThrMax}{\ensuremath{3.50\times10^{-4}}}
\renewcommand{\nBonfThr}{\ensuremath{1.67\times10^{-4}}}
\renewcommand{\nFwerRatioMin}{0.48}
\renewcommand{\nFwerRatioMax}{1.13}
\renewcommand{\nMissMvnZero}{0.0501}
\renewcommand{\nMissWyZero}{0.0625}
\renewcommand{\nMissMvnHalf}{0.0877}
\renewcommand{\nMissWyHalf}{0.1475}
\renewcommand{\nMissMvnFull}{0.1168}
\renewcommand{\nMissWyFull}{0.2200}
\renewcommand{\nMissMvnStrFull}{0.1658}
\renewcommand{\nMissWyStrFull}{0.2825}
\renewcommand{\nMissBonfMod}{0.1850}
\renewcommand{\nMissBonfStr}{0.2475}
\renewcommand{\nMissRatioMvnHalf}{1.8}
\renewcommand{\nMissRatioWyHalf}{2.9}
\renewcommand{\nMissLamLo}{0.926}
\renewcommand{\nMissLamHi}{1.141}
\renewcommand{\nMissCrossover}{75}
\renewcommand{\nCmpACATInd}{0.0526}
\renewcommand{\nCmpACATStr}{0.0547}
\renewcommand{\nCmpACATCal}{0.0518}
\renewcommand{\nCmpFisherInd}{0.0501}
\renewcommand{\nCmpFisherStr}{0.1566}
\renewcommand{\nCmpFisherCal}{0.0372}
\renewcommand{\nCmpStoufferInd}{0.0496}
\renewcommand{\nCmpStoufferStr}{0.1805}
\renewcommand{\nCmpStoufferCal}{0.0323}
\renewcommand{\nCmpminPInd}{0.0523}
\renewcommand{\nCmpminPStr}{0.0452}
\renewcommand{\nCmpminPCal}{0.0537}
\renewcommand{\nCmpTFisherInd}{0.0438}
\renewcommand{\nCmpTFisherStr}{0.1214}
\renewcommand{\nCmpTFisherCal}{0.0461}
\renewcommand{\nCmpOmnibusInd}{0.0404}
\renewcommand{\nCmpOmnibusStr}{0.1371}
\renewcommand{\nCmpOmnibusCal}{0.0342}
\renewcommand{\nCmpWorst}{Stouffer}
\renewcommand{\nCmpBest}{minP}
\renewcommand{\nCmpTFisherNA}{4.0}
\renewcommand{\nCandN}{18,563}
\renewcommand{\nCandNGwas}{18,563}
\renewcommand{\nCandNPath}{9,727}
\renewcommand{\nCandCtrlFound}{13}
\renewcommand{\nCandCtrlTotal}{13}
\renewcommand{\nCandCtrlMedian}{20}
\renewcommand{\nCandCtrlTopPct}{12}
\renewcommand{\nCandRandom}{9,282}
\renewcommand{\nCandLDLR}{1}
\renewcommand{\nCandPCSKnine}{3}
\renewcommand{\nCandGwasTies}{0}
}{}

%% \maybefig{file}{width}{caption}{label} -- draws the figure if the PNG has
%% been generated, otherwise a visible reminder box, so the manuscript
%% compiles before the simulations have been run.
\newcommand{\maybefig}[4]{%
  \IfFileExists{../figures/#1}{\includegraphics[width=#2]{#1}}%
  {\fbox{\parbox[c][0.28\textheight][c]{#2}{\centering\textcolor{red}%
   {\textbf{Figure pending:} \texttt{\detokenize{#1}}\\[2pt]
    generated by \texttt{scripts/05\_figures.R}}}}}%
  \caption{#3}\label{#4}}

%% \maybeinput{file} -- \input an auto-generated table if it exists.
\newcommand{\maybeinput}[1]{\IfFileExists{#1}{\input{#1}}{%
  \begin{table}[!t]\centering\textcolor{red}{\textbf{Table pending:}
  \texttt{\detokenize{#1}}}\end{table}}}

\begin{document}

\journaltitle{Bioinformatics}
\DOI{DOI added during production}
\copyrightyear{2026}
\pubyear{2026}
\vol{XX}
\issue{x}
\access{Published: Date added during production}
\appnotes{Original Paper}
\firstpage{1}

\title[CATFISH with Westfall--Young calibration]{Family-wise error control for
correlated pathway tests in GWAS: CATFISH with Westfall--Young calibration}

\author[1,3,$\ast$]{Nirwan Tandukar}
\author[2]{Jung-Ying Tzeng}
\author[3]{Rub\'en Rell\'an-\'Alvarez\ORCID{0000-0000-0000-0000}}

\address[1]{\orgdiv{Department of Molecular and Structural Biochemistry}, \orgname{North Carolina State University}, \orgaddress{\street{Raleigh}, \postcode{27606}, \state{NC}, \country{USA}}}
\address[2]{\orgdiv{Department of Statistics}, \orgname{North Carolina State University}, \orgaddress{\street{Raleigh}, \postcode{27606}, \state{NC}, \country{USA}}}
\address[3]{\orgdiv{Genetics and Genomics Program}, \orgname{North Carolina State University}, \orgaddress{\street{Raleigh}, \postcode{27606}, \state{NC}, \country{USA}}}

\corresp[$\ast$]{Nirwan Tandukar. \href{mailto:ntanduk@ncsu.edu}{ntanduk@ncsu.edu}}
\corresp[$\ast$]{Rub\'en Rell\'an-\'Alvarez. \href{mailto:rrellan@ncsu.edu}{rrellan@ncsu.edu}}

\received{Date}{0}{Year}
\revised{Date}{0}{Year}
\accepted{Date}{0}{Year}

\abstract{
\textbf{Motivation:} Pathway analysis of genome-wide association study (GWAS)
summary statistics must contend with two distinct sources of dependence. Genes
within a pathway are correlated through linkage disequilibrium (LD), and
pathways are correlated with one another through shared gene membership. Most
pipelines address only the first, calibrating each pathway against its own
LD-aware null and then applying a Bonferroni or Benjamini--Hochberg correction
that assumes the resulting pathway $p$-values are independent. They are not,
and the correction is consequently miscalibrated in a direction that depends on
how strongly the pathway family overlaps.\\
\textbf{Results:} We extend CATFISH, a correlation-aware omnibus pathway
framework combining ACAT, Fisher, adaptive soft TFisher, minP and Stouffer
statistics, with a Westfall--Young-style family-wise calibration. A single
genome-wide correlated null draw, generated from a multivariate normal model of
gene--gene correlation, is propagated through every pathway in the family; read
down its columns it gives each pathway a calibrated marginal $p$-value, and read
across its rows the minimum over pathways gives a family-wise threshold that
accounts for between-pathway dependence without modelling it. Numerical
precision is treated as a correctness requirement: in a GWAS of $n\approx1.3$
million individuals the omnibus statistic underflows double precision for the
majority of pathways, and the resulting ties silently destroy the ranking.
Across simulated null families spanning three LD regimes and three pathway
sizes, independence-based calibration was anti-conservative while per-pathway
and Westfall--Young calibration held type~I and family-wise error near nominal,
and both remained stable when up to four of the five component tests were
deliberately miscalibrated. Under five stylised signal architectures, the
component predicted to be sensitive to each was the best-performing component in
$\nArchCellsIn$ of $\nArchCellsTot$ architecture-by-size combinations. Applied to a maize
grain-nitrogen GWAS, the Westfall--Young ordering recovered L-homoserine
biosynthesis as the leading pathway and placed a known aspartate kinase isozyme
first among 34,873 genes. Applied to a GLGC LDL-cholesterol GWAS, it exposed a
regime in which the self-contained null is false almost everywhere: 11.2\% of
genes reach genome-wide significance, essentially every pathway is significant,
and exact FWER control does not rescue the ranking. We argue this is a property
of self-contained testing at very large $n$, not a failure of multiple-testing
correction.\\
\textbf{Availability:} CATFISH is an R package available at
\url{https://github.com/nirwantandukar/CATFISH} under the MIT licence. All
simulation, analysis and figure code for this paper is in the
\texttt{catfish2/} directory of that repository.\\
\textbf{Contact:} \href{mailto:ntanduk@ncsu.edu}{ntanduk@ncsu.edu}\\
\textbf{Supplementary information:} Supplementary data are available at
\textit{Bioinformatics} online.}

\keywords{genome-wide association study, pathway analysis, gene-set test,
family-wise error rate, Westfall--Young, permutation, ACAT, MAGMA}
\keywords[Abbreviations]{GWAS, MAGMA, LD, FWER, WY, MVN, ACAT}

\boxedtext{Key Messages}{
\begin{itemize}
\item Pathway $p$-values are correlated because pathways share genes;
Bonferroni and Benjamini--Hochberg corrections applied to them are therefore
miscalibrated, and the direction of the error depends on the overlap structure
of the particular pathway database used.
\item A Westfall--Young max-statistic null, built by propagating one
genome-wide correlated draw through the whole pathway family, controls the
family-wise error rate in the strong sense without assuming the pathways are
independent; its validity rests instead on the multivariate normal correlation
model used to generate the draw.
\item Implementing the omnibus in log space is a correctness requirement, not an
optimisation: at GWAS sample sizes above roughly one million the analytic
omnibus underflows double precision and the pathway ranking collapses into ties.
\item Self-contained pathway tests saturate at very large $n$. When 11\% of all
genes are genome-wide significant, a random 50-gene set contains several of
them, almost every pathway rejects, and exact FWER control does not restore
discrimination. Saturation is a property of the hypothesis being tested, not of
the correction applied to it.
\end{itemize}}

\maketitle

%%%%%%%%%%%%%%%%%%%%%%%%%%%%%%%%%%%%%%%%%%%%%%%%%%%%%%%%%%%%%%%%%%%%%%%
\section{1 Introduction}

Genome-wide association studies (GWAS) have transformed complex trait research
by identifying thousands of loci linked to disease, development, and
quantitative phenotypes. Yet most variants explain only a small fraction of
phenotypic variance and often fall in noncoding regions, limiting single-marker
mechanistic interpretation \citep{Manolio2009,Boyle2017}. Complex traits are
often highly polygenic, with relevant signals distributed across many
modest-effect loci rather than concentrated in a few genome-wide significant
variants \citep{Akula2025,Lappalainen2024}. Gene-set analysis, including
pathway-based approaches, is therefore a key post-GWAS strategy for translating
summary statistics into biological hypotheses. By aggregating evidence across
predefined groups of functionally related genes, pathway-based methods enhance
interpretability, reduce the multiple-testing burden, and reveal coordinated
pathway disruptions that single-locus analyses may miss
\citep{Mooney2015,deLeeuw2015}. This is especially valuable for polygenic
traits, where many genes may influence the same phenotype through shared
biological processes. In such settings, a candidate gene becomes more compelling
when locus- or gene-level association evidence is reinforced by pathway-level
enrichment from a biologically relevant process. Pathway support therefore
provides functional context for prioritization, although it is best interpreted
as complementary rather than definitive evidence because pathway annotations
remain incomplete and often overlapping.

A plethora of pathway analysis tools has been developed for GWAS. Classical
over-representation methods test whether predefined pathways are enriched among
significant genes, as implemented in DAVID, g:Profiler, and WebGestalt
\citep{Sherman2022,Kolberg2023,Liao2019}. Ranking-based methods such as Gene Set
Enrichment Analysis (GSEA) assess whether pathway genes are non-randomly
concentrated at the top or bottom of a ranked gene list
\citep{Subramanian2005}. GWAS-focused extensions, including MAGENTA, INRICH, and
MAGMA, add features such as SNP-to-gene mapping, pathway-level modeling, and
adjustment for linkage disequilibrium (LD) or gene size
\citep{Segre2010,Lee2012,deLeeuw2015}. Accounting for LD is essential in GWAS
pathway analysis because nearby variants are often correlated and therefore do
not provide independent evidence of association. If LD is ignored, clusters of
linked SNPs can be effectively counted multiple times, leading to inflated
gene-level aggregation and enrichment biases toward larger genes or regions with
dense marker coverage. Several widely used post-GWAS methods therefore
explicitly account for LD through model-based or reference-based dependence
estimates rather than assuming independence. VEGAS uses multivariate normal
simulation, Pascal uses LD from a reference population to compute gene and
pathway scores, and MAGMA models gene-level dependence through an approximated
correlation structure
\citep{Liu2010,deLeeuw2015,Lamparter2016,Watanabe2017}.

These methods address dependence \emph{within} a pathway. A second source of
dependence operates \emph{between} pathways and is far less often handled.
Curated pathway databases are built from overlapping annotations: a gene that
participates in several metabolic steps is a member of several pathways, and
hierarchical resources such as Reactome place parent and child pathways in the
same family. The pathway-level $p$-values returned by any method are therefore
correlated with one another, sometimes strongly. Nonetheless, the final
multiple-testing step in essentially every published pipeline is a Bonferroni or
Benjamini--Hochberg adjustment of those $p$-values, both of which are derived
under independence or under positive-regression dependence conditions that
overlapping gene sets do not generally satisfy. The resulting family-wise error
rate is not the nominal one, and the size of the discrepancy is a property of
the particular database rather than of the method, which makes it invisible to
the kind of simulation study that generates independent pathways.

The classical solution to exactly this problem is the Westfall--Young
max-statistic procedure \citep{Westfall1993}. Rather than correcting each
$p$-value in isolation, one generates a null realization of the \emph{entire}
family of tests, records the single most extreme statistic observed in that
realization, and repeats. The distribution of that maximum is the reference
distribution for the observed minimum, and the resulting threshold controls the
family-wise error rate under the complete null whatever the dependence structure
among the tests, because the dependence is inherited from the way the null
realizations are generated rather than assumed. In a GWAS pathway setting this means drawing
correlated gene-level null statistics genome-wide, propagating that single draw
through every pathway in the family, and recording the minimum pathway
$p$-value. Nothing about the between-pathway correlation needs to be estimated,
modelled, or even named.

Here we extend CATFISH (Combining \textbf{C}auchy combination test (ACAT),
\textbf{A}daptive \textbf{T}Fisher (soft) test, \textbf{F}isher's test,
m\textbf{I}n-P, and \textbf{S}touffer's method for \textbf{H}olistic pathway
analysis), a correlation-aware pathway analysis framework that integrates
multiple complementary gene-set statistics in a unified omnibus strategy, with
Westfall--Young family-wise calibration. Different pathway tests are sensitive
to different genetic architectures: some are best at detecting broad,
distributed enrichment, whereas others are more sensitive to sparse strong
signals or heterogeneous evidence. Instead of relying on a single summary
statistic, CATFISH evaluates each pathway with several complementary tests and
combines them in an omnibus framework. This improves robustness when the true
signal configuration is unknown while preserving the interpretability of the
component tests. Beyond pathway-level inference, CATFISH prioritizes candidate
genes by integrating locus-level GWAS signals, gene-level MAGMA results, and
pathway-level support.

This paper makes three contributions beyond the original CATFISH framework.
First, we add the Westfall--Young layer described above and characterize how it
behaves: the same correlated null, read down its columns, supplies each
pathway's marginal $p$-value, and read across its rows, the family-wise one. Second, we show that the entire statistic must be computed in log
space. This sounds like an implementation detail and is not: at the sample sizes
now routine in human GWAS consortia, the analytic omnibus $p$-value for a
typical pathway falls below the smallest representable double-precision number,
every such pathway is assigned the same floor value, and the ranking that the
whole analysis exists to produce is destroyed without any error being raised.
Third, we report a negative result we believe is general. In a GLGC
LDL-cholesterol GWAS of roughly 1.3 million individuals, the self-contained null
hypothesis is false for almost every gene set that could be constructed, nearly
every pathway is significant, and exact family-wise error control does not
recover a usable ordering. We interpret this as saturation of the self-contained
hypothesis rather than as a defect of the correction, and discuss what it
implies for the choice between self-contained and competitive testing as GWAS
sample sizes continue to grow.

CATFISH is implemented as an R package and is freely available at
\url{https://github.com/nirwantandukar/CATFISH} under the MIT licence.

\subsection*{1.1 Pathway signal archetypes}

CATFISH was developed to capture heterogeneous pathway-level signal patterns,
ranging from highly sparse gene-level associations to weak, broadly distributed
enrichment. We define five conceptual pathway signal archetypes to clarify how
different gene-level profiles are expected to favor different component tests.
These archetypes are not intended as empirically learned classes of real
pathways but as statistical templates for interpreting test behavior and
motivating the complementary design of CATFISH.

For interpretive clarity, we consider a pathway-level working null and
alternative. Under the null, gene-level evidence within a pathway matches the
background distribution after accounting for dependence, with no systematic
excess of small $p$-values. Under the alternative, the pathway shows enrichment
through one or more forms of excess association evidence, including extreme
individual genes, broad moderate signals, weak diffuse shifts, or mixed
patterns. Figure~\ref{fig:archetypes} illustrates the five ranked signal shapes,
and Table~\ref{tab:archetype_summary} summarizes their statistical profiles and
expected CATFISH component sensitivity.

In short, Strong Sparse Signal (SSS) refers to pathways driven by a few strongly
associated genes that favor extreme-tail methods. Single-Gene Proxy (SGP), a
special case of SSS, is composed of a signal dominated by one leading gene, but
it should be interpreted cautiously as gene-centric. Moderate Dense Signal (MDS)
has many modest, consistent gene associations that favor aggregation methods.
Weak Diffuse Enrichment (WDE) constitutes subtle pathway-wide shifts, with few
individually significant genes, favoring distributed-signal methods. Mixed
Signal Architecture (MSA) represents one or a few strong signals plus broader
moderate support that favor balanced methods.

\begin{figure*}[!t]
\centering
\maybefig{Fig0_archetypes.png}{\textwidth}{Ranked gene-level association
profiles for five pathway signal archetypes. Panels illustrate conceptual
within-pathway $-\log_{10}(p)$ profiles for Strong Sparse Signal (SSS),
Moderate Dense Signal (MDS), Weak Diffuse Enrichment (WDE), Mixed Signal
Architecture (MSA), and Single-Gene Proxy (SGP). Horizontal lines mark
representative significance thresholds. These archetypes motivate the use of
complementary CATFISH component tests because different methods are
preferentially sensitive to different ranked signal shapes.}{fig:archetypes}
\end{figure*}

\begin{table*}[!t]
\centering
\caption{Summary of pathway signal archetypes and expected CATFISH component
sensitivity. Each archetype represents a conceptual ranked gene-level
association profile rather than an empirically learned class of real pathways.}
\label{tab:archetype_summary}
\scriptsize
\setlength{\tabcolsep}{4pt}
\begin{tabular}{p{0.17\textwidth} p{0.30\textwidth} p{0.21\textwidth} p{0.13\textwidth}}
\hline
\textbf{Archetype} & \textbf{Gene-level $p$-value pattern} & \textbf{Ranked-profile shape} & \textbf{Expected sensitive tests} \\
\hline
Strong Sparse Signal (SSS) &
A small number of genes have extremely small $p$-values, while most genes remain close to null. &
Sharp early peak followed by a rapid drop; pronounced elbow. &
ACAT, minP \\
Moderate Dense Signal (MDS) &
Many genes show modest but consistent association, without one clearly dominant gene. &
Broad shoulder of moderate signals across a substantial fraction of the pathway. &
Fisher, Stouffer, soft TFisher \\
Weak Diffuse Enrichment (WDE) &
Most genes are individually non-significant, but the pathway shows a subtle overall shift toward smaller $p$-values. &
Gentle pathway-wide deflection with few or no strong individual peaks. &
Stouffer, Fisher \\
Mixed Signal Architecture (MSA) &
One or a few strong genes occur together with a broader set of moderate supporting signals. &
Strong leading peak plus a moderate shoulder. &
soft TFisher, Fisher, ACAT \\
Single-Gene Proxy (SGP) &
One dominant gene accounts for most of the pathway-level signal, while the remaining genes are near null. &
Single extreme peak followed by an otherwise null-like profile. &
minP, ACAT \\
\hline
\end{tabular}
\end{table*}

%%%%%%%%%%%%%%%%%%%%%%%%%%%%%%%%%%%%%%%%%%%%%%%%%%%%%%%%%%%%%%%%%%%%%%%
\section{2 Methods}

\subsection*{2.1 Gene-level inputs and component tests}

For a pathway $S$ containing $m=|S|$ genes, let $p_g$ denote the gene-level
association $p$-value for gene $g\in S$ and $Z_g$ the corresponding gene-level
$Z$-statistic. We write $\mathbf{p}_S=(p_1,\dots,p_m)$ and
$\mathbf{Z}_S=(Z_1,\dots,Z_m)$. Gene-level statistics are obtained with MAGMA
\citep{deLeeuw2015}, which aggregates SNP-level GWAS evidence across variants
assigned to a gene while accounting for LD, gene size, and local marker density,
and which additionally returns the gene--gene correlation matrix induced by LD.
These quantities are the common inputs for every CATFISH component test and for
both calibration layers.

Each pathway is assessed using five complementary component tests (Table~\ref{tab:components}): ACAT \citep{Liu2019ACATV,Liu2020ACAT}, Fisher \citep{Fisher1932}, adaptive soft-thresholding TFisher \citep{Zhang2020TFisher}, minP, and Stouffer \citep{Stouffer1949}. Together, they cover the archetypes in Section~1.1 and apply directly to gene-level $p$-values or $Z$-scores without external weights or pathway-specific tuning. We combine the component $p$-values into a single omnibus statistic via an ACAT-type cross-method operator $\mathcal{O}(\cdot)$. This omnibus is a pathway-level summary of evidence across components, not an additional gene-level test.

Four components (ACAT, Fisher, minP, Stouffer) are fully determined by gene-level $p$-values and require no tuning. Soft TFisher is the exception: it adapts over a grid of truncation thresholds $\tau$, which determines which genes contribute. This is the only tunable choice in CATFISH, and we justify it below.

Following \citet{Zhang2020TFisher}, the grid spans small, intermediate and large thresholds rather than concentrating at one end: we use \(\tau\in\{10^{-2},\,5\times10^{-2},\,5\times10^{-1}\}\) and take the minimum calibrated TFisher \(p\)-value over the grid. The original authors' grid also includes \(\tau=1\), which we omit because \(\tau=1\) is Fisher's method and Fisher is already carried as a separate component. Spanning the range matters because the number of genes surviving truncation is \(\mathrm{Binomial}(m,\tau)\): a grid confined to small \(\tau\) leaves \(m\tau\ll1\) at realistic set sizes, the statistic is then zero on a large fraction of pathways, and the component is degenerate precisely where it is most often applied.

A threshold is informative mainly when \(m\tau\) is not far below one, so which \(\tau\) values are active depends on set size: \(\tau=5\times10^{-1}\) is informative from the smallest sets upward, whereas \(\tau=10^{-2}\) requires \(m\approx100\). Adapting across \(\tau\) rather than fixing a single value is what allows one component to serve both sparse and diffuse architectures, and users can assess the grid for their own database using \(m\tau\).

\smallskip\noindent\textit{Undefined components.} When no gene in a pathway falls below any threshold in the grid, the soft TFisher statistic is zero at every \(\tau\) and its one-sided \(p\)-value is exactly one. We treat this as \emph{undefined} rather than as evidence against enrichment, and the omnibus for that pathway is formed from the components that are defined. The distinction is not cosmetic. The ACAT operator maps each component \(p\)-value through \(\tan\{(1/2-p)\pi\}\), so a value at the upper boundary contributes a large negative term; a component that could not be computed would otherwise vote against the pathway and bias the omnibus toward the null. The rule is applied per pathway rather than by set size, because a large pathway in which no gene is individually strong is equally uncomputable. It does condition the component on having at least one gene below a threshold, so where TFisher does contribute its \(p\)-values are stochastically small; because the resampling null is generated under the identical rule, that conditioning appears in the null distribution and is absorbed by the calibration. An analytic omnibus \(p\)-value under this rule would be anti-conservative, whereas the resampled one is not.

\smallskip\noindent\textit{Precision limit.} The soft TFisher implementation returns the null distribution function of the soft-thresholding statistic, so the one-sided \(p\)-value is formed as one minus that value. In double precision the difference is exactly zero once the distribution function exceeds \(1-10^{-16}\), and a strongly enriched pathway would then report \(p=0\). We cap the soft TFisher \(p\)-value at \(10^{-16}\), the smallest value the computation can resolve, rather than at the general lower floor applied to components whose tails are computed in closed form. A cap below the precision limit asserts resolution the statistic does not have; Section~4.3 shows what that does to the family-wise threshold.

\begin{table*}[!t]
\centering
\caption{CATFISH pathway-level component statistics. For a pathway $S$ with
$m=|S|$ genes and gene-level $p$-values $p_g$, each component emphasizes a
different pathway signal pattern. The analytic reference distributions define
the statistics; final inference uses the dependence-aware calibrations of
Section 2.3.}
\label{tab:components}
\small
\renewcommand{\arraystretch}{1.25}
\begin{tabular}{p{0.14\textwidth} p{0.33\textwidth} p{0.23\textwidth} p{0.16\textwidth}}
\hline
\textbf{Component} & \textbf{Pathway statistic} & \textbf{Analytic reference} & \textbf{Primary sensitivity} \\
\hline
ACAT &
\[ T_{\mathrm{ACAT}}(S)=\frac{1}{m}\sum_{g\in S}\tan\!\left\{\pi\left(\tfrac12-p_g\right)\right\} \]
&
\[ p_{\mathrm{ACAT}}(S)=\tfrac12-\frac{1}{\pi}\arctan T_{\mathrm{ACAT}}(S) \]
& Sparse or extreme-tail signals. \\
Fisher &
\[ T_{\mathrm{F}}(S)=-2\sum_{g\in S}\log p_g \]
& Under independence, $T_{\mathrm{F}}\sim\chi^2_{2m}$. &
Broad enrichment with many small-to-moderate $p$-values. \\
Soft TFisher &
\[ W(S;\tau)=\sum_{g\in S}\left[-2\log p_g+2\log\tau\right]_{+} \]
& Adaptive over a $\tau$ grid; minimum calibrated $p$-value is taken. &
Intermediate signals concentrated in a subset of genes. \\
Stouffer &
\[ Z_{\mathrm{St}}(S)=\frac{1}{\sqrt m}\sum_{g\in S}Z_g,\quad Z_g=\Phi^{-1}(1-p_g) \]
& $p_{\mathrm{St}}(S)=1-\Phi\!\left(Z_{\mathrm{St}}(S)\right)$. &
Weak, diffuse, pathway-wide enrichment. \\
minP &
\[ T_{\min}(S)=\min_{g\in S}p_g \]
& Under independence, $p_{\min}=1-(1-T_{\min})^m$. &
Single strongest gene or highly sparse signals. \\
\hline
\end{tabular}
\end{table*}

\subsection*{2.2 Two sources of dependence}

CATFISH statistics are correlated at two levels. The two arise from different
causes, and both have to be addressed before any pathway $p$-value means what
it claims to.

\emph{Within a pathway}, component tests are dependent for two reasons: each is a deterministic function of the same $\mathbf{p}_S$, inducing cross-method dependence even if genes were independent; and gene-level evidence is correlated via LD, overlapping SNP-to-gene assignments, and shared local genomic architecture. Thus, any analytic reference distribution assuming independence is wrong for every component and typically anti-conservative, because correlated genes contain less information than the same number of independent genes.

\emph{Between pathways}, pathway-level $p$-values are correlated because databases assign genes to multiple sets. A gene shared across pathways propagates its (possibly extreme) null draw to all of them, shifting their evidence in the same direction. The strength and pattern of this dependence depend on database structure: hierarchies with nested sets create widespread but weak dependence, whereas largely disjoint curated collections create rare but strong dependence among overlapping sets. Consequently, independence-based multiple-testing corrections (Bonferroni or Benjamini--Hochberg under its positive-dependence conditions) need not achieve nominal error rates; the discrepancy is driven by the database, not the method.

These are often treated as separate problems, but in CATFISH both reduce to the same object: a vector of correlated gene-level statistics. A single genome-wide null draw, carrying the estimated LD, is scored through all pathways at once. Within a pathway, it provides null components and accounts for within-pathway dependence; across pathways, it provides a null for the family minimum and accounts for between-pathway dependence. No correlation structure is estimated or modeled; both are inherited. Section~2.3 develops this.

The construction relies on three assumptions: (i) \emph{Approximate normality}: under the null, the pathway-level vector satisfies $\mathbf{Z}_S\sim\mathcal{N}(0,R_S)$. (ii) \emph{Correlation-matrix adequacy}: correlation matrices from MAGMA-reported gene pairs provide an LD-informed working approximation to the unobserved gene-level covariance; unreported pairs are set to zero, and adequacy is evaluated via calibration and robustness analyses below. (iii) \emph{No residual structural confounding}: gene-level inputs are adjusted so pathway signals are not primarily driven by gene size or SNP density.

\subsection*{2.3 Calibration and family-wise control from a shared correlated null}

CATFISH obtains two different $p$-values for every pathway: a marginal one,
which asks whether that pathway is enriched, and a family-wise one, which asks
whether evidence this strong would appear anywhere among the $K$ pathways
tested by chance. Both are read from a single set of resampled null draws. The
procedure has four steps.

\paragraph{Step 1: generate the null.}
Gene--gene correlations are read from MAGMA's \texttt{.genes.raw} output as a
list of gene pairs and their correlations. For each chromosome $c$ containing
$m_c$ genes we form an $m_c\times m_c$ correlation matrix $R_c$ by starting
from the identity and filling in the reported pairs,
\[
R_c = I_{m_c} + \sum_{(i,j)\in\mathcal{P}_c}
      r_{ij}\left(e_ie_j^{\top}+e_je_i^{\top}\right),
\]
where $\mathcal{P}_c$ is the set of reported pairs for chromosome $c$; pairs not in $\mathcal{P}_c$ remain zero. We compute correlations only for gene pairs close enough to share LD, so $R_c$ is banded by construction: zeros for distant genes are approximations, not imputations, whereas zeros for nearby but unreported pairs introduce approximation error (see Section~5). Genes on different chromosomes are treated as independent, which is structurally correct. If $R_c$ is not positive definite, we floor eigenvalues at $10^{-6}$, reconstruct, and renormalise the diagonal; the implementation reports affected chromosomes.

Writing $R_c=U_c^{\top}U_c$ for the Cholesky factor, replicate $b$ draws
\[
\mathbf{Z}^{(b)}_c = U_c^{\top}\boldsymbol{\varepsilon},\qquad
\boldsymbol{\varepsilon}\sim\mathcal{N}(0,I_{m_c}),
\]
independently per chromosome, and the concatenation $\mathbf{Z}^{(b)}$ is one
genome-wide null realisation carrying the estimated LD structure. Gene-level
null $p$-values follow by the Gaussian copula
$p_g^{(b)}=1-\Phi(Z_g^{(b)})$, matching MAGMA's one-sided sign convention. This
is repeated for $b=1,\dots,B$.

We draw genome-wide, not pathway-by-pathway, because Step~4 requires all $K$ pathways to be scored on the \emph{same} realisation so shared genes perturb them jointly. Drawing each pathway from its own correlation submatrix would be marginally correct but would remove the dependence structure Step~4 needs.

\paragraph{Step 2: score every pathway on every replicate.}
For each replicate we recompute all five component statistics and the omnibus
operator $\mathcal{O}(\cdot)$ for every pathway in the family
$\mathcal{F}=\{S_1,\dots,S_K\}$, giving a $B\times K$ matrix of null omnibus
values
\[
Q^{(b)}_k = \mathcal{O}\!\left(\mathbf{p}^{(b)}(S_k)\right),
\qquad b=1,\dots,B,\quad k=1,\dots,K .
\]
Because all components are recomputed from the same draw, $Q$ preserves within-pathway gene dependence, cross-method dependence, and—since all $K$ pathways share the draw—between-pathway dependence. The observed values $q^{\mathrm{obs}}_k=\mathcal{O}(\mathbf{p}^{\mathrm{obs}}(S_k))$ are computed the same way from the real data. Applying $\mathcal{O}$ to the observed components without resampling yields the uncalibrated analytic omnibus $\mathrm{Omnibus}_{\mathrm{An}}(S_k)=q^{\mathrm{obs}}_k$, which assumes gene independence and serves as the baseline.

\paragraph{Step 3: the marginal \texorpdfstring{$p$}{p}-value.}
Compare each pathway's observed statistic with its own null values across the
$B$ replicates. The LD-aware marginal $p$-value is the Monte Carlo tail
probability
\[
\tilde p_k \;=\;
\frac{1+\sum_{b=1}^{B}\mathbf{1}\!\left(Q^{(b)}_k\le q^{\mathrm{obs}}_k\right)}{B+1}.
\]
This quantity is computed from the genome-wide draw rather than from the pathway-specific submatrix $R_{S_k}$; because $R_{S_k}$ is simply the $R_c$ submatrix for that pathway's genes, both use the same correlation information, so nothing is lost. We call $\tilde p_k$ the per-pathway calibrated $p$-value and use it to rank pathways. It corresponds to what is often called MVN calibration; we avoid that label because the multivariate normal model is the generator of the null draws, not a separate calibration step, and the same draws serve Step~4.

\paragraph{Step 4: the family-wise \texorpdfstring{$p$}{p}-value.}
Within each replicate, record the smallest statistic across the whole family,
\[
M^{(b)} = \min_{1\le k\le K} Q^{(b)}_k ,
\]
and compare each observed statistic with the resulting distribution. This gives
the Westfall--Young adjusted $p$-value \citep{Westfall1993}
\[
\tilde p^{\mathrm{WY}}_k \;=\;
\frac{1+\sum_{b=1}^{B}\mathbf{1}\!\left(M^{(b)}\le q^{\mathrm{obs}}_k\right)}{B+1},
\]
Rejecting all pathways with $\tilde p^{\mathrm{WY}}_k\le\alpha$ controls the family-wise error rate at level $\alpha$ under the complete null, whatever the dependence among pathways, because a false rejection anywhere in the family requires the family minimum to cross the threshold. Three further points apply. First, the control is in the strong sense, not only under the complete null. Westfall--Young procedures control the family-wise error rate strongly when the joint distribution of the statistics for the true null hypotheses does not depend on which other hypotheses are false---the condition of subset pivotality \citep{Westfall1993}. Two features of the construction give it. Each pathway statistic is a deterministic function of its own genes' $z$-scores only; and under the model of Step~1, with association represented as a shift in the means of the associated genes, the genes of a pathway containing no associated gene are jointly $\mathcal{N}(0,R)$, restricted to those genes, whatever happens elsewhere. The observed statistics of the null pathways therefore have exactly their complete-null joint distribution, and because the minimum over any subset of pathways is never smaller than the minimum over the whole family, the probability that any null pathway crosses the threshold is at most $\alpha$. The argument requires the soft TFisher truncation grid to be fixed in advance: a grid chosen from the empirical distribution of all gene-level $p$-values, as in data-adaptive implementations, would make every pathway's statistic depend on genes outside it. Second, the null is parametric: it is generated from the multivariate normal model of Step~1 rather than by permuting phenotypes, so the procedure is Westfall--Young in construction but inherits the adequacy of that model. Third, the counting form of the adjusted $p$-value above, in which null minima equal to the observed value count as at least as extreme, is required rather than a matter of taste. A threshold taken as a quantile of $M$ and applied with an inclusive comparison rejects every observation tied at that quantile, and if the omnibus loses resolution in the tail so that many null minima coincide, the empirical family-wise error rate can approach one (Section~4.3).

Between-pathway correlation is never explicitly estimated or modelled. Because all $K$ pathways are evaluated on the same draw in Step~2, genes shared across pathways perturb them jointly, as in the observed data, and this dependence is reflected in the distribution of $M$. The threshold $\hat t_\alpha=\mathrm{quantile}_\alpha(M)$ is the family-wise significance cutoff implied by the pathway database in use; its ratio to the Bonferroni cutoff $\alpha/K$ quantifies the database’s departure from independence at no extra cost, provided it is read on calibrated statistics; on the analytic scale the ratio also absorbs any miscalibration of the per-pathway statistic (Section~4.1).

Steps~3 and~4 use the same matrix $Q$ and require no additional resampling. Their outputs relate as a raw and adjusted $p$-value: $\tilde p_k$ assesses pathway $S_k$ alone, while $\tilde p^{\mathrm{WY}}_k$ accounts for testing $K$ pathways. Both are reported; neither replaces the other.

\paragraph{Size standardisation.}
A minimum across tests is meaningful only when the tests share a scale. The raw omnibus does not: Fisher’s statistic increases with $m$, so the Step~4 family minimum would mechanically favor the largest sets. We therefore also report a variant that first standardizes each observed value within its own column of $Q$ as a $z$-score, $z_k=(q^{\mathrm{obs}}_k-\mu_k)/\sigma_k$, where $\mu_k$ and $\sigma_k$ are that column’s mean and standard deviation, and then takes the row minimum. This rescales values for the minimum operation; it is not a pathway-size bias correction, and Section~4.1 shows it does not behave as one. In the simulations of Section~2.6 every pathway in a configuration has the same size and dependence structure, so the pathways are exchangeable, share one marginal null, and the row minimum is taken on the raw scale; standardisation is needed only when pathway sizes differ, as in real data. We standardise by location and scale rather than by each pathway's empirical null distribution function. The latter would put pathways on an exactly common scale, but would resolve each to only $1/(B+1)$, which is coarser than the family minimum requires and would reintroduce ties.

\subsection*{2.4 Numerical implementation in log space}

The omnibus statistic must be computed entirely in log space. This is a correctness requirement, not an optimisation, and it prevents silent failure.

The smallest positive normalized double is $\approx2.2\times10^{-308}$. In a GWAS with $n\approx1.3$ million, gene-level $p<10^{-300}$ is common: in the GLGC LDL-cholesterol data (Section~3.9), the published $p$-value is exactly zero for 473 chromosome~1 variants, including the \textit{SORT1} lead rs12740374 at $-\log_{10}p=1725$. Combining even a few such genes with Fisher’s method yields a pathway statistic whose double-precision $p$-value underflows to zero. All affected pathways then share the same floor value, ties are broken by storage order, and the ranking the analysis aims to produce is silently corrupted. In our first implementation, the top-ranked pathways were (on inspection) alphabetical.

Five changes are needed. (i) Compute all component statistics on $\ell_g=\log p_g$ and return log $p$-values via \texttt{stats::pchisq(\dots,\,log.p=TRUE)} and \texttt{stats::qnorm(\dots,\,log.p=TRUE)}. (ii) Evaluate the ACAT sum using the small-$p$ asymptote: for $p<10^{-15}$, use $\tan\{\pi(\tfrac12-p)\}\approx 1/(p\pi)$, accumulating those terms as a log-sum-exp of $-\ell_g-\log\pi$; use the trigonometric form for the rest. (iii) Replace $p=\tfrac12-\arctan(T)/\pi$ with $p\approx1/(T\pi)$ for $T>10^{6}$, where cancellation near $\tfrac12$ dominates truncation error. (iv) Since ACAT can be negative when most inputs exceed $0.5$, handle that case in a separate branch computing $\tfrac12-\arctan(T)/\pi$ directly; clamping would set $p=0.5$ and destroy ordering in the upper half. (v) The minP inversion $1-(1-q)^m$ underflows for $q\lesssim10^{-16}/m$, causing minP to drop out; use the $m q$ asymptote whenever $\log m+\ell_{\min}<\log 10^{-8}$.

These changes let the pipeline represent pathway $p$-values down to $e^{-745}\approx10^{-324}$ and, crucially, preserve their \emph{ordering} well below where probability-scale computations fail. They are implemented in the scripts used for the real-data analyses. The simulations of Section~2.6 instead call the component functions of the CATFISH package directly, so that they validate the code a user runs; those functions work on the probability scale with a lower floor of $10^{-300}$, which is adequate there because no simulated null approaches it. Soft TFisher is subject to its own, much higher, precision limit (Section~2.1).

\subsection*{2.5 Candidate-gene prioritization}

The output a biologist acts on is a short list of genes, so CATFISH ranks genes
by a composite score over three evidence layers:
\[
\mathrm{Score}(g)=S_{\mathrm{GWAS}}(g)+S_{\mathrm{MAGMA}}(g)
+0.5\,S_{\mathrm{PATH}}(g)+0.3\,B_{\mathrm{hub}}(g).
\]
Each layer enters by rank rather than by its raw statistic,
$S_\bullet(g)=-\log_{10}\!\left(r_\bullet(g)/N_\bullet\right)$, where
$r_{\mathrm{GWAS}}$ is the gene's rank by minimum SNP-level $p$-value,
$r_{\mathrm{MAGMA}}$ its rank by gene-level $p$-value, and $r_{\mathrm{PATH}}$
the rank of its best pathway; rank one is the strongest evidence. The hub term
$B_{\mathrm{hub}}(g)=\log_{10}(n_{\mathrm{path}}(g)+1)$ rewards membership in
several pathways. Genes missing a layer score zero for it, so an absent layer
neither helps nor penalises.

We use ranks instead of $-\log_{10}p$ to make layers with different tail behaviour comparable and to avoid the $p$-value floors in Section~2.4. Pathway evidence is halved because the pathway $p$-value is a set-level statistic shared across member genes; the hub term is intentionally secondary. We fixed weights \emph{a priori}; simulations showed rankings were stable for $w_{\mathrm{path}}\in[0,1]$ and $w_{\mathrm{hub}}\in[0,0.5]$, with defaults on a broad plateau rather than a fragile optimum \citep{Tandukar2025Thesis}. We did not re-tune on the validation genes in Section~3.9.4 to avoid circularity.

The score is a heuristic ranking metric, not an inferential model: it has no null distribution or $p$-value. We also report per-layer ranks, because the GWAS and MAGMA layers use the same variants and are not independent evidence.

\subsection*{2.6 Simulation design}

Simulations are run on the omnibus statistic under its null, except where a signal is injected deliberately (item~6). Two dependence designs are used. In the \emph{shared-universe} design, pathways are overlapping subsets of one gene universe: $G=3000$ genes in independent blocks of 50 with exchangeable within-block correlation $\rho$, standing in for chromosomes and LD blocks, from the first 600 of which we draw $K=150$ pathways of $m$ genes each at random. Each gene therefore belongs to several pathways, a pathway spans several blocks, and a single genome-wide draw moves every pathway containing a given gene---the between-pathway dependence the family-wise procedure is designed to control. In the \emph{exchangeable} design, used only in item~1, each pathway is drawn independently and all $m$ of its genes are mutually correlated at $\rho$. That design has no between-pathway dependence; it isolates within-pathway correlation at an intensity well beyond what block-structured LD produces, and serves as a stress test. Each study reports two summaries of the null matrix: per-pathway marginal type~I error (by columns) and family-wise error rate (by rows), with Bonferroni as the family-wise comparator. Thresholds are fitted on one set of draws and evaluated on a held-out set, so rates are not correct by construction. Component statistics are computed with the CATFISH package functions, using the soft TFisher grid, undefined-component rule and precision cap of Section~2.1. We do not simulate genotypes or phenotypes or run a full GWAS; that is future work (Section~5) and is not needed for the questions here. Each study corresponds to a script in \texttt{catfish2/scripts/} and writes a CSV to \texttt{catfish2/results/} from which the tables and figures are generated.

\begin{enumerate}
\item \textbf{Null calibration and type~I error} (\texttt{01\_null\_calibration.R}), exchangeable design. Under the global null we compare four decision rules: the analytic independence-based $p$-value; per-pathway calibration against a null bank shared by all pathways of a given size; Bonferroni; and the Westfall--Young adjusted $p$-value of Step~4, with the null distribution of the family minimum estimated on one half of the simulated families and evaluated on the other. We vary $\rho\in\{0,0.3,0.6\}$ and $m\in\{5,20,50\}$ and report type~I error at $\alpha\in\{0.05,0.01\}$, family-wise error, the genomic-control inflation factor $\lambda=\mathrm{median}(\chi^2_1\text{-transformed }p)/0.4549$, and the number of distinct values taken by the family minimum, which detects loss of resolution directly.

\item \textbf{Component-breakage stress test} (\texttt{02\_breakage\_stress.R}).
We miscalibrate $k$ of five components by replacing each affected component’s $p$-value with $p^{\gamma}$, and test whether the omnibus and its calibration remain valid. We run both corruption directions: for $0<p<1$, $\gamma<1$ yields $p^{\gamma}>p$ (conservative) and $\gamma>1$ yields $p^{\gamma}<p$ (anti-conservative), the latter being most relevant for false positives. We use $\gamma\in\{0.7,1.4\}$; sweep $k=0,\dots,4$; evaluate up to three component combinations per $k$; and cross with the same LD regimes at $m\in\{20,50\}$. We apply corruption to both null replicates and observed statistics, allowing resampling calibration to absorb it.

\item \textbf{Leave-one-out component sensitivity} (\texttt{03\_adaptive\_loo.R}).
We remove each component in turn, recompute the null bank and the calibration
from scratch, and measure the change in type~I error and $\lambda$. This
isolates how much each component contributes to, or destabilises, the calibrated
omnibus, and identifies components whose removal would change inference.

\item \textbf{Per-component calibration} (\texttt{09\_component\_calibration.R}).
Rather than combining, we calibrate each of the five components against its own
null and report analytic and calibrated type~I error, family-wise rate and
$\lambda$ per component. This identifies which statistics LD actually breaks and
whether a single calibrated component is usable in place of the omnibus when the
signal architecture is known in advance.

\item \textbf{Robustness to an incomplete correlation file}
(\texttt{08\_missing\_correlation.R}). Data are simulated under the true correlation, but calibration (null bank, marginal $p$-values, and Westfall--Young threshold) uses a matrix where a fraction $\{0,0.25,0.5,0.75,1\}$ of off-diagonal entries is set to zero, mimicking unreported gene pairs. Evaluation draws are shared across missingness levels; only calibration varies. For real data, \texttt{10\_pair\_coverage.R} places a correlation file on this scale by reporting the fraction of within-pathway gene pairs on the same chromosome and close enough to be correlated for which MAGMA reports a value.

\item \textbf{Archetype power} (\texttt{17\_archetype\_power\_wy.R}), shared-universe design with $\rho=0.3$. One pathway in each family carries signal shaped as one of the five archetypes of Section~1.1; the remaining $K-1$ are null. Signal is a mean shift on the gene-level $z$ scale applied to a fixed fraction of the pathway's genes, at positions permuted in every replicate: SSS shifts $15\%$ of genes by $3.0$; MDS $60\%$ by $1.1$; WDE $90\%$ by $0.55$; MSA $10\%$ by $1.8$ and a further $60\%$ by $1.1$; and SGP $5\%$ by $4.0$, with at least one gene in every layer. Shifts are multiplied by an effect scale over a grid running from undetectable to saturated, at $m\in\{5,20,50\}$. For each component and for the omnibus we record the probability that the signal pathway clears its own Westfall--Young threshold. Regret is the difference in that probability between the best single component and the omnibus in the same cell.
\end{enumerate}

Unless stated otherwise, each configuration uses $B$ null draws for the
calibration bank and a separate set of replicate datasets for evaluation, with
both values set at the top of each script so that a reduced smoke-test
configuration can be run with \texttt{bash run\_all.sh quick}.

\subsection*{2.7 Real-data analyses}

\textbf{Maize grain nitrogen.} Gene-level statistics were computed with MAGMA
from a mixed-linear-model (MLM) GWAS of grain nitrogen content, using a
$10\,\mathrm{kb}$ upstream and downstream window for SNP-to-gene assignment.
Pathways were taken from CornCyc \citep{Hawkins2021CornCyc,Caspi2008BioCyc},
giving $172$ pathways after size filtering. Gene--gene correlations were
recovered per chromosome from MAGMA's \texttt{.genes.raw} output; we verified
explicitly that all ten chromosomes were represented, after an earlier run in
which only chromosome~1 correlations had been parsed and $85\%$ of the genome
was silently treated as independent. Candidate-gene scoring used the three-layer
score of Section~2.5 over all $34{,}873$ annotated genes.

\textbf{Human LDL cholesterol.} We used the Global Lipids Genetics Consortium
LDL-cholesterol summary statistics for the European-ancestry meta-analysis
($n\approx1.3$ million) \citep{Graham2021GLGC}.

\emph{Variant filtering and the $p$-value floor.} Variants were restricted to
rsIDs present in the 1000 Genomes European reference panel used for LD. The
released \texttt{pvalue} column underflows to exactly zero at the strongest
loci, so all probabilities were reconstructed from \texttt{pvalue\_neg\_log10}.
For the MAGMA gene analysis, whose \texttt{--pval} input requires a
probability, these were floored at $10^{-300}$; MAGMA is unaffected because the
floor is far beyond any gene-level resolution it reports. The candidate-gene
GWAS layer is affected, because it ranks genes by their minimum variant
$p$-value and the floor collapses the strongest genes into a single tie. That
layer therefore reads an unfloored $-\log_{10}p$ file written by
\texttt{11\_munge\_neglog10.sh}; Section~3.9 reports what the floor costs when
it is not removed.

\emph{Gene analysis.} Variants were assigned to genes with a $10\,\mathrm{kb}$
upstream and downstream window using the NCBI37.3 gene locations. MAGMA gene
analysis was run per chromosome with \texttt{--batch} and combined with
\texttt{--merge}, which preserves the gene--gene correlations in the merged
\texttt{.genes.raw} file; running the chromosomes as independent jobs and
concatenating the outputs would discard them.

\emph{Pathways.} Pathway inference used the Reactome subset of MSigDB C2:CP
($1{,}390$ sets after size filtering), with HumanCyc additionally converted into
the CATFISH pathway format for comparison. We report Reactome because a
hierarchical database is where between-pathway dependence is strongest, which is
the regime the family-wise procedure is designed for.

\emph{Candidate genes.} The GWAS layer uses MAGMA's own
\texttt{.genes.annot} file rather than re-deriving a window, so that the GWAS
and gene layers are assigned from identical variant sets. The pathway layer is
ranked on the observed omnibus statistic rather than the Westfall--Young
adjusted $p$-value. The two give the identical ordering, since the adjusted
$p$-value is a monotone transform of the statistic, but the adjusted value is
floored at $1/(B+1)$ and the entire top of the list is tied there at the $B$
used here.

\emph{Positive controls.} Thirteen genes with established roles in LDL-C
biology---\textit{LDLR}, \textit{PCSK9}, \textit{APOB}, \textit{APOE},
\textit{SORT1}, \textit{HMGCR}, \textit{NPC1L1}, \textit{ABCG8}, \textit{LPA},
\textit{CETP}, \textit{LDLRAP1}, \textit{MYLIP} and \textit{ANGPTL3}---were
named before the prioritization was run and are used to assess it. They are a
validation set, not a discovery; the score has no access to them.

%%%%%%%%%%%%%%%%%%%%%%%%%%%%%%%%%%%%%%%%%%%%%%%%%%%%%%%%%%%%%%%%%%%%%%%
\section{3 Results}

\subsection*{3.1 Independence-based calibration is anti-conservative;
per-pathway and Westfall--Young calibration are not}

Under the global null, the analytic omnibus $p$-value is valid only if the genes
within a pathway are independent. Table~\ref{tab:null_calibration} and
Figure~\ref{fig:type1} show what happens when they are not. At $\rho=0$ and
$m=5$ the analytic type~I error was $\nTAnInd$ against a nominal $0.05$, which
is the behaviour one expects when the generating assumption holds. As either the
within-pathway correlation or the pathway size increases, the analytic
calibration departs from nominal, reaching $\nTAnStr$ at $\rho=0.6$, $m=50$,
the largest value in the grid. The direction is the expected one:
correlated genes carry less independent information than the analytic reference
distribution assumes, so the statistic looks more extreme than it is.

Per-pathway calibration against the correlated null removed most of this.
Empirical type~I error was $\nTMvnInd$ at $\rho=0$, $m=5$ and $\nTMvnStr$ at
$\rho=0.6$, $m=50$, remaining within $[\nTMvnMin,\nTMvnMax]$ across the whole
grid. At the stricter level $\alpha=0.01$, the analytic calibration reached
$\nTAnMaxOne$ against the calibrated statistic's $\nTMvnMaxOne$, indicating that the
discrepancy grows rather than shrinks in the tail, which is where pathway
analysis operates.

The family-wise columns test the Westfall--Young construction, and here they are
a genuine test rather than a check of arithmetic: the null distribution of the
family minimum is estimated on half of the simulated families and applied to the
other half. Empirical family-wise error ranged from $\nTWyMin$ to $\nTWyMax$
across the grid, within the Monte Carlo error of $250$ held-out families, for
which the standard error at the nominal level is about $0.014$. Bonferroni, applied to the analytic $p$-value, inherits that
value's miscalibration: its family-wise error was $\nTBonfInd$ in the
independent configuration with $m=5$, where the analytic $p$-value is valid, but
$\nTBonfMod$ at $\rho=0.3$, $m=20$ and $\nTBonfStr$ at $\rho=0.6$, $m=50$. The
pathways in this design are independent of one another, so none of that
inflation comes from between-pathway dependence; all of it is the per-test
statistic being wrong, compounded across the family.

Reaching these numbers required two corrections worth recording, because both
failures were silent. In an earlier version the family-wise error at
$\rho=0.6$, $m=50$ was $1.00$. The cause was not the Westfall--Young logic but a
loss of resolution in the omnibus. Soft TFisher's tail underflowed for about
$3.5\%$ of null pathways in that configuration, each such pathway was assigned
the same floor value, and with $200$ pathways per family nearly every family's
minimum was that one shared number; a quantile of a distribution concentrated on
a single value, applied with an inclusive comparison, rejects every family.
Capping soft TFisher at the precision it actually has (Section~2.1) and
computing the adjusted $p$-value in its counting form (Section~2.3) restored the
family minimum to a continuous statistic: even in the worst configuration it
takes $\nWyDistinctMin$ distinct values across the $500$ simulated families.

\maybeinput{auto/tab_null.tex}

\begin{figure*}[!t]
\centering
\maybefig{Fig1_type1_error.png}{\textwidth}{Empirical type~I error of the
CATFISH omnibus under the global null at $\alpha=0.05$, by within-pathway LD
level ($\rho=0$, $0.3$, $0.6$) and pathway size ($m=5$, $20$, $50$). The
horizontal line marks the nominal level. The analytic independence-based
calibration departs from nominal as either correlation or pathway size
increases; per-pathway calibration against the correlated null, and the
Westfall--Young family-wise error rate, both track the nominal
level.}{fig:type1}
\end{figure*}

\subsection*{3.2 Genomic-control inflation tracks the same failure}

The genomic-control inflation factor $\lambda$ summarises the same null in a
form familiar from GWAS practice (Figure~\ref{fig:lambda}), and it points the
wrong way. For the analytic omnibus, $\lambda$ was $\nLamAnInd$ at $\rho=0$,
$m=50$ and fell to $\nLamAnStr$ at $\rho=0.6$, $m=50$---far \emph{below}
one---in exactly the configuration where the analytic type~I error was largest.
The two summaries are not in conflict. $\lambda$ is computed from the median of
the null $p$-values and type~I error from their lower tail, and under strong
within-pathway correlation the analytic null distribution is not shifted but
reshaped: excess mass in the lower tail together with a median above one half.
An analyst reading $\lambda$ alone would conclude the test was conservative.
After per-pathway calibration $\lambda$ lay in $[\nLamMvnMin,\nLamMvnMax]$. A
pathway analysis reporting $\lambda$ near or below one is therefore not evidence
that its tail is calibrated.

\begin{figure}[!t]
\centering
\maybefig{Fig2_lambda.png}{\columnwidth}{Genomic-control inflation factor
$\lambda$ for the analytic and per-pathway calibrated omnibus under the global null, by
LD level and pathway size. The horizontal line marks $\lambda=1$.}{fig:lambda}
\end{figure}

\subsection*{3.3 The omnibus degrades gracefully when its components are broken}

A multi-test omnibus is only useful if a single misbehaving component cannot
dominate it, and a family-wise procedure is only useful if that robustness
survives at the family level too. We tested both by corrupting components on
purpose, replacing the $p$-value of each of $k$ selected components with
$p^{\gamma}$ and running the conservative ($\gamma=0.7$) and the
anti-conservative ($\gamma=1.4$) direction over the overlapping pathway universe
of Section~2.6, in which each gene belongs to $\nBrkMembership$ pathways on
average (Table~\ref{tab:breakage}, Figure~\ref{fig:breakage}).

With no components corrupted ($k=0$), the analytic omnibus had mean marginal
type~I error $\nBrkAnKzero$ and the calibrated omnibus $\nBrkMvnKzero$; at the
family level the Westfall--Young rate was $\nBrkWyKzero$ and Bonferroni
$\nBrkBonfKzero$. Under anti-conservative corruption of four of the five
components ($k=4$, $\gamma=1.4$) the analytic omnibus degraded to
$\nBrkAnKfour$, while the calibrated marginal rate held at $\nBrkMvnKfour$ and
the Westfall--Young family-wise rate at $\nBrkWyKfour$ against Bonferroni's
$\nBrkBonfKfour$. Under conservative corruption the analytic omnibus moved the
other way, to $\nBrkConsAnKfour$, losing power rather than control, while the
calibrated marginal and family-wise rates held at $\nBrkConsMvnKfour$ and
$\nBrkConsWyKfour$. Across both directions and every configuration the
family-wise rate stayed within $[\nBrkWyRangeLo,\nBrkWyRangeHi]$ and the
marginal rate within $[\nBrkMvnRangeLo,\nBrkMvnRangeHi]$. Genomic-control
inflation for the same runs (Figure~\ref{fig:breakage_lambda}) tracks the
analytic statistic's drift in both directions while staying near one for the
calibrated statistic.

The interpretation is specific and worth stating carefully. Resampling
calibration is robust to component miscalibration because the \emph{same}
corruption is applied to the null replicates as to the observed data, so a
component that is uniformly miscalibrated shifts the null and the observation
together and largely cancels; the family-wise threshold, being a quantile of
the same corrupted null, moves with them. This is a real and useful property,
but it protects against systematic miscalibration, not against a component that
is wrong in a way that correlates with the signal of interest. A component whose
error depended on the data would not cancel, and we have not tested that case.
The analytic calibration has no such protection and inherits the corruption
directly, in whichever direction it was applied.

\maybeinput{auto/tab_breakage.tex}

\begin{figure*}[!t]
\centering
\maybefig{Fig3_breakage.png}{\textwidth}{Component-breakage stress test on
overlapping pathways. Rows are the direction of the deliberate corruption
($p\mapsto p^{\gamma}$, with $\gamma=1.4$ anti-conservative and $\gamma=0.7$
conservative); columns are the within-pathway LD level. Each panel shows the
per-pathway marginal rates and the family-wise rates together, since all four
are rates at the nominal $0.05$ marked by the dotted line. Line type
distinguishes pathway size. The analytic statistic and the Bonferroni
correction applied to it both drift with $k$; the calibrated marginal rate and
the Westfall--Young family-wise rate do not.}{fig:breakage}
\end{figure*}

\begin{figure}[!t]
\centering
\maybefig{Fig3b_breakage_lambda.png}{\columnwidth}{Genomic-control inflation
$\lambda$ for the component-breakage runs, by corruption direction and LD
level. $\lambda$ is a property of the marginal $p$-value distribution and has
no family-wise counterpart, so it is shown separately from
Figure~\ref{fig:breakage}. The analytic statistic drifts with $k$ in whichever
direction the components were corrupted; the calibrated statistic stays near
one.}{fig:breakage_lambda}
\end{figure}

\subsection*{3.4 No single component is load-bearing}

Removing each component in turn and rebuilding the null bank, the marginal
calibration and the family-wise threshold from scratch
(Table~\ref{tab:loo}, Figure~\ref{fig:loo}) left the marginal type~I error
within $[\nLooMvnMin,\nLooMvnMax]$ against $\nLooFullMvn$ for the full
five-component omnibus, and the Westfall--Young family-wise rate within
$[\nLooWyMin,\nLooWyMax]$ against $\nLooFullWy$. The largest departure from
nominal followed removal of \texttt{\nLooWorst} ($\nLooWorstVal$).

The same table carries the sharpest contrast in the paper between the two
family-wise procedures. Averaged over pathway sizes and leave-one-out variants,
the Westfall--Young rate was $\nLooWyInd$, $\nLooWyMod$ and $\nLooWyStr$ in the
independent, moderate and strong LD regimes---flat at nominal---while the
Bonferroni rate moved from $\nLooBonfInd$ to $\nLooBonfMod$ to $\nLooBonfStr$,
reaching $\nLooBonfMax$ at worst. We are careful about what this does and does
not show. Bonferroni is applied here to the analytic omnibus $p$-value, which is
itself anti-conservative under LD (Section~3.1), so the inflation is the
compounding of an uncalibrated per-test statistic with a correction that cannot
repair it, which is what a pipeline that skips calibration actually does.
In the independent rows, where the analytic $p$-value is valid, Bonferroni was
close to nominal: overlap at an average of $\nBrkMembership$ pathways per gene
is not enough to make it noticeably conservative. The lesson is that a family-wise correction inherits
whatever miscalibration the per-test statistic carries, whereas the
Westfall--Young threshold, being a quantile of the same null the statistic came
from, does not.

Within that pattern, removing Stouffer, Fisher or soft TFisher lowered the
Bonferroni rate, while removing ACAT or minP raised it. Stouffer and Fisher are
the aggregation statistics that within-pathway correlation inflates
(Section~3.5), and soft TFisher's uncalibrated value is a minimum over thresholds
and therefore stochastically small; ACAT and minP stay close to nominal, so
removing them leaves the omnibus more dominated by the inflated components.

No single deletion broke error control at either level, which is the expected
consequence of calibrating empirically: the null bank is recomputed for
whichever components remain, so both the marginal $p$-value and the family-wise
threshold remain valid for any subset. What changes with the subset is
sensitivity, not validity. The practical implication is that users may drop a
component they distrust for a particular dataset---for example minP when
gene-level $p$-values have been floored at a reporting limit---without
invalidating the inference, at some cost in power against the architectures that
component is best at.

Genomic-control inflation behaves the same way: the fourth panel of
Figure~\ref{fig:loo} shows $\lambda$ for the calibrated marginal statistic
staying near one for every deletion, so no component is load-bearing for the
inflation diagnostic either.

\maybeinput{auto/tab_loo.tex}

\begin{figure*}[!t]
\centering
\maybefig{Fig4_leave_one_out.png}{\textwidth}{Leave-one-out component
sensitivity at $m=50$, with each component removed in turn and the null bank,
marginal calibration and Westfall--Young threshold rebuilt from scratch. From
left: marginal type~I error, Westfall--Young family-wise error rate, Bonferroni
family-wise error rate, and genomic-control inflation $\lambda$ for the
calibrated marginal statistic. The dashed line marks the nominal level ($0.05$,
or $\lambda=1$). Note the y-scale of the third panel, which is the point of the
comparison.}{fig:loo}
\end{figure*}

\subsection*{3.5 Linkage disequilibrium breaks the aggregation tests, not the
extreme-value tests}

The omnibus hides which of its inputs is failing. Calibrating each component
separately (Table~\ref{tab:components_cal}, Figure~\ref{fig:components}) shows
that the damage LD does is concentrated in two of the five, and that the
identity of those two is predictable from what the statistics measure.

Under the independent null every component was correctly sized (for example
ACAT at $\nCmpACATInd$ and Stouffer at $\nCmpStoufferInd$), which is the
expected behaviour when the generating assumption holds. Under strong LD the
two aggregation statistics---Fisher, which sums $-2\log p$ across the pathway,
and Stouffer, which sums $Z$-scores---inflated to $\nCmpFisherStr$ and
$\nCmpStoufferStr$ against the nominal $0.05$. The two extreme-value
statistics barely moved: ACAT went to $\nCmpACATStr$ and minP to
$\nCmpminPStr$. The mechanism is transparent. An aggregation statistic adds one
term per gene and implicitly treats each as independent evidence, so correlated
genes are counted repeatedly and the statistic grows; a statistic that depends
on the most extreme gene is insensitive to how many near-duplicates of it the
pathway contains. The omnibus inherited the inflation of its worst inputs,
reaching $\nCmpOmnibusStr$.

Calibration against the correlated null removed this for every component,
bringing Fisher to $\nCmpFisherCal$, Stouffer to $\nCmpStoufferCal$ and the
omnibus to $\nCmpOmnibusCal$. The practical consequence is that a user who knows a pathway's signal
architecture may use a single calibrated component in place of the omnibus
without losing error control, which is what makes the architecture-specific
testing described in Section~2.3 usable.

This result also explains an observation in Section~3.4: removing an
aggregation component lowered the Bonferroni family-wise rate while removing
ACAT or minP raised it. The aggregation components are the ones that were
inflated in the first place.

Three implementation notes belong here, all about soft TFisher. First, the
statistic is undefined on a pathway in which no gene falls below any truncation
threshold. Under the grid of Section~2.1 this did not occur at $m\ge20$ in any
configuration, but it is frequent in small pathways under grids confined to
small $\tau$, and a single undefined draw in a calibration bank makes that
pathway's calibrated $p$-value undefined unless non-finite draws are filtered
before the tail count, which the implementation does. Second, an undefined
component must be excluded from the omnibus rather than entered at $p=1$
(Section~2.1). Because ACAT weights by $\tan\{(1/2-p)\pi\}$, entering the
boundary value makes a test that could not be computed act as evidence against
the pathway; under the null it pushes the median omnibus $p$-value above one
half, so that $\lambda$ falls below one while the tail is simultaneously
inflated. Third, soft TFisher's uncalibrated value is a minimum over thresholds
rather than a $p$-value, so its analytic columns in
Table~\ref{tab:components_cal} are not interpretable as calibration; only its
calibrated and family-wise columns are.

\maybeinput{auto/tab_components_cal.tex}

\begin{figure*}[!t]
\centering
\maybefig{Fig7_component_calibration.png}{\textwidth}{Per-component
calibration at $m=50$. Left: type~I error at $\alpha=0.05$ under the analytic
independence assumption. Centre: the same components calibrated against the
correlated null. Right: genomic-control inflation $\lambda$ for the analytic
statistics. The aggregation components, Fisher and Stouffer, are the ones LD
inflates; the extreme-value components, ACAT and minP, are nearly
unaffected. Calibration restores all five.}{fig:components}
\end{figure*}

\subsection*{3.6 Both calibrations degrade when the correlation file is
incomplete, and the family-wise one degrades faster}

Every result above assumes the gene--gene correlation matrix is the one the
data were generated from. In practice it is assembled by filling MAGMA's
reported gene pairs into an identity matrix, so a pair MAGMA does not report is
silently treated as independent. We measured the cost directly by generating
data from the true correlation structure while calibrating from a matrix with a
known fraction of its off-diagonal entries zeroed
(Table~\ref{tab:missing}, Figure~\ref{fig:missing}).

With a complete matrix the marginal type~I error was $\nMissMvnZero$ and the
Westfall--Young family-wise rate $\nMissWyZero$, both at nominal. Dropping half
the pairs raised them to $\nMissMvnHalf$ and $\nMissWyHalf$; dropping all of
them, which is the independence assumption in its pure form, gave
$\nMissMvnFull$ and $\nMissWyFull$ under moderate LD and $\nMissMvnStrFull$ and
$\nMissWyStrFull$ under strong LD.

Expressed as multiples of each procedure's own nominal level---the only fair
comparison between a per-pathway rate and a per-family one---the marginal
calibration was $\nMissRatioMvnHalf\times$ nominal at $50\%$ missingness and
the family-wise procedure $\nMissRatioWyHalf\times$. Westfall--Young is
therefore more sensitive to an incomplete correlation file than the marginal
calibration it is built from. The reason is structural rather than
incidental: the family minimum over $K$ pathways is a tail statistic, so it
depends on the extreme of the null distribution, which is precisely where an
error in the correlation structure does the most damage. The marginal
$p$-value depends only on the bulk.

This does not overturn the case for family-wise control, because the comparison
that matters is against the correction it replaces rather than against the
marginal rate it is not competing with. Bonferroni's family-wise rate in the
same runs was $\nMissBonfMod$ and $\nMissBonfStr$ and does not depend on the
correlation file at all, since it is computed from the analytic $p$-value.
Westfall--Young remained the better family-wise procedure up to $\nMissCrossover\%$
missingness, and lost only when essentially no correlation information
remained---at which point it is calibrating against a null that does not
resemble the data.

It does, however, change how the procedure should be reported. A
Westfall--Young $p$-value is only as trustworthy as the correlation file behind
it, and the completeness of that file is measurable rather than assumed. The
quantity to check is narrow: a correlation matrix from MAGMA is \emph{supposed}
to be sparse, because correlations are computed only for gene pairs close
enough to share LD, so a zero between distant genes is correct rather than
missing. What matters is the fraction of within-pathway gene pairs that are on
the same chromosome and close enough to be correlated for which a value is
actually reported. We provide \texttt{10\_pair\_coverage.R} for this, and
recommend that it be run and its result stated alongside any family-wise
$p$-value.

Finally, and consistently with Section~3.2, $\lambda$ did not flag any of this.
Across the missingness sweep it ranged from $\nMissLamLo$ to $\nMissLamHi$,
drifting \emph{below} one as correlation information was removed, while the
marginal type~I error more than doubled. This is the second independent setting
in which $\lambda$ at or below one coexists with an inflated tail, and it is the
clearest argument we can make that $\lambda$ should not be used as the sole
calibration diagnostic in pathway analysis.

\maybeinput{auto/tab_missing.tex}

\begin{figure*}[!t]
\centering
\maybefig{Fig6_missing_correlation.png}{\textwidth}{Robustness to an incomplete
gene--gene correlation file. Data are generated from the true correlation
structure while calibration uses a matrix with the stated fraction of
off-diagonal entries zeroed. Left and centre: error rates at $\alpha=0.05$
under moderate and strong LD. Right: genomic-control inflation $\lambda$ for
the marginal calibration, which falls below one rather than rising as the
type~I error more than doubles.}{fig:missing}
\end{figure*}

\subsection*{3.7 Each archetype is detected by the component predicted for it}

Section~1.1 proposed five pathway signal archetypes and predicted, on
statistical grounds, which component tests should be sensitive to each. That
prediction is testable only under simulation, because the architecture of a real
pathway is unknown. We generated pathways carrying each archetype's ranked
signal shape, embedded them in a family of null pathways drawn from the same
correlated gene universe, and measured the probability that the signal pathway
clears the family-wise Westfall--Young threshold. Effect sizes were scaled over
a grid so that every archetype passes from undetectable to saturated, and the
component set, the universe and the resampling are exactly those used
elsewhere in this paper.

The predictions hold. In $\nArchCellsIn$ of $\nArchCellsTot$ archetype-by-size
cells the component that most often achieved the highest family-wise power was
inside the set predicted in Table~\ref{tab:archetype_summary}
(Table~\ref{tab:archverify}). The separation is sharp and in the expected
direction: diffuse enrichment is recovered by Stouffer, with the extreme-tail
tests near their null; single-gene and strongly sparse signals are recovered by
minP and ACAT, with Fisher and Stouffer close to powerless; and dense or mixed
architectures by Fisher, with soft TFisher close behind. No component is best
everywhere, which is the premise the omnibus rests on.

The ranks below first place matter as much as the winners. Across the
$\nArchInfCells$ informative cells, Fisher finished first $\nFisherFirst$ times
but fourth or lower in $\nFisherFourthPlus$, and minP finished last in
$\nMinPLast$. Soft TFisher finished first only $\nTFFirst$ time, but second in
$\nTFSecond$ and third in $\nTFThird$, and never lower than third: it is the
only component never among the weakest. Inside an omnibus that averages its
inputs, that is the more useful property. A component that is never best but
never near its null protects the combination precisely in the cells where the
specialists disagree. It also qualifies Table~\ref{tab:archetype_summary}: soft
TFisher belongs among the tests sensitive to dense and mixed signal, where it
ran close behind Fisher, but it was not the most sensitive test for any
archetype.

What matters in practice is not which component wins but how much the omnibus
gives up by not knowing in advance which one would. We measure that as
\emph{regret}: the difference between the best single component in a cell and
the omnibus in the same cell. Regret is small for most architectures
(Table~\ref{tab:archregret}, Figure~\ref{fig:archregret}), from
$\nArchRegretBest$ for $\nArchRegretBestA$ to $\nArchRegretWorstA$ at
$\nArchRegretWorst$. Against the alternative of choosing one component in
advance, which risks the full gap between the best and worst component in a
cell---frequently more than $0.5$ in power---that is a favourable trade.

The exception is worth stating plainly. In the single-gene case the omnibus can
give up as much as $\nArchMaxRegret$ in a single cell, because minP and ACAT
detect the pathway while Fisher and Stouffer are near their null, and averaging
across components dilutes a signal that only two of them can see. This is a
property of combination rather than a defect of any component, and it is the
price of not having to choose. It also marks the regime in which a gene-level
test, rather than a pathway-level one, is the appropriate instrument.

\maybeinput{auto/tab_archetype_verification.tex}
\maybeinput{auto/tab_archetype_regret.tex}

\begin{figure*}[!t]
\centering
\maybefig{Fig11_archetype_power.png}{\textwidth}{Family-wise power to detect a
signal pathway, by archetype (columns) and pathway size (rows), as a function of
effect scale. Each line is one component test; the dashed black line is the
ACAT-O omnibus. The component predicted for each archetype in
Table~\ref{tab:archetype_summary} separates from the others over the range in
which power is informative.}{fig:archpower}
\end{figure*}

\begin{figure}[!t]
\centering
\maybefig{Fig12_archetype_regret.png}{\columnwidth}{Power the omnibus gives up
relative to the best single component in each cell. Bars are the mean across
sizes and effect scales; crosses mark the worst single cell. The dashed line
marks a regret of $0.05$.}{fig:archregret}
\end{figure}


\subsection*{3.8 Maize grain nitrogen: Westfall--Young recovers the known
pathway and gene}

We applied the full pipeline to a mixed-linear-model GWAS of maize grain
nitrogen content, testing $172$ CornCyc pathways. The leading pathway under the
size-standardized Westfall--Young ordering was \texttt{HOMOSERSYN-PWY},
L-homoserine biosynthesis (13 genes; omnibus $p=1.58\times10^{-3}$;
size-standardized family-wise $p=0.070$), followed by \texttt{PWY-5097},
L-lysine biosynthesis VI (25 genes; omnibus $p=3.04\times10^{-3}$;
size-standardized family-wise $p=0.338$).

This ordering reproduces, independently and from a different statistical
framework, the result of a previously reported analysis of the same trait
\citep{Tandukar2025Nitrogen}, in which L-homoserine biosynthesis was the only
pathway enriched across all populations examined and L-lysine biosynthesis VI
was the runner-up for reasons that analysis sets out. Recovering the same two
pathways in the same order out of $172$ candidates, using a method that shares
no calibration machinery with the original, is the strongest external check we
can presently make on the pipeline.

The three-layer candidate score of Section~2.5, computed over all $34{,}873$
annotated maize genes with pathway evidence taken from the Westfall--Young
adjusted $p$-values, placed \texttt{Zm00001eb220040} first. This gene is the
aspartate kinase isozyme named in the earlier analysis, and L-homoserine
biosynthesis was its highest-scoring pathway, so the gene- and pathway-level
layers agree rather than merely coexisting.

Two caveats are essential. First, nothing in this analysis is significant: the
leading pathway's size-standardized family-wise $p$-value is $0.070$, and under
the unstandardized Westfall--Young ordering it is $0.284$ against a Bonferroni
value of $0.272$. The result is a ranking that agrees with prior biology, not a
discovery. Second, grain nitrogen is highly polygenic and the MLM GWAS contained
no genome-wide significant genes at all (Table~\ref{tab:density}), which is
precisely the regime in which a ranking is the only available output. We regard
the agreement as encouraging evidence that the ordering is meaningful in the
low-signal regime, and we do not claim more than that.

\subsection*{3.9 Human blood lipids: saturation at the pathway level, intact
resolution at the gene level}

The GLGC lipid meta-analyses ($n\approx1.3$ million) are the opposite regime
from the maize analysis: very large samples, well-characterised traits, and
gene-level signal everywhere. We analysed LDL-C and HDL-C through an identical
pipeline---same reference panel, same gene annotation, same Reactome sets, same
window, same number of permutations---so that the only thing differing between
them is the phenotype. Together they produce one clearly negative result and
one clearly positive one, and the pair says more than either alone.

\subsubsection*{3.9.1 The self-contained null is false almost everywhere}

At the gene level, $44.1\%$ of tested genes had $p<0.05$ for LDL-C and $11.2\%$
reached genome-wide significance; for HDL-C the figures were $52.7\%$ and
$14.8\%$, with a median gene-level $p$-value of $0.038$. The consequence for
self-contained pathway testing is arithmetical rather than statistical. A
\emph{randomly chosen} set of fifty genes contains on average $5.6$ genome-wide
significant genes under LDL-C and $7.4$ under HDL-C, and contains at least one
with probability above $0.99$ in both. The self-contained null
hypothesis---that the genes in a set carry no more association than the
background---is therefore false for essentially any set of genes one could
assemble, including sets with no biological coherence whatever.

The output reflects this. Of the $1{,}390$ Reactome pathways, $1{,}366$ were
significant under Westfall--Young family-wise control at $\alpha=0.05$ for
LDL-C and $1{,}384$ for HDL-C. The natural first explanation is an inadequate
multiple-testing correction, and we can rule it out: the Westfall--Young
procedure controls the family-wise error rate exactly under the realized
between-pathway dependence, and it removed essentially nothing. The rejections
are not false positives in the sense a correction addresses. The null really is
false for those sets. The problem is the hypothesis, not the threshold.

A second, subtler consequence is that the ranking acquires a strong dependence
on pathway size. The twenty highest-ranked pathways had a median of $159$ genes
for LDL-C and $182$ for HDL-C, against a median of $29$ across all $1{,}390$
sets, and the correlation between $\log_{10}$ set size and the pathway
statistic was $-0.825$. In a saturated genome a large set samples more of an
enriched background, so it accumulates more evidence for a reason that has
nothing to do with the trait. Any interpretation of the ordering has to be made
against that gradient, and we return to it below.

\subsubsection*{3.9.2 Two traits, different genes, the same pathway}

LDL-C and HDL-C are biologically distinct arms of lipoprotein metabolism, and
their gene-level results differ accordingly
(Figure~\ref{fig:convergence}A). Across the $18{,}563$ genes analysed in both,
the Spearman correlation of gene-level $Z$ statistics was $0.393$. Of the top
hundred genes in each trait, $13$ were shared, and of the genes reaching
genome-wide significance in either, the Jaccard overlap was $0.208$
($830$ shared of $3{,}987$). These are not two readings of the same signal.

Despite that, the highest-ranked pathway of the $1{,}390$ tested was the same
in both: \texttt{REACTOME\_PLASMA\_LIPOPROTEIN\_ASSEMBLY\_REMODELING\_AND\_%
CLEARANCE}. At $74$ genes it is well below the $159$--$182$ median size of the
surrounding top twenty, so it reaches first place against the size gradient of
Section~3.9.1 rather than because of it.

The reason the two traits converge is visible in which genes carry the pathway
in each (Figure~\ref{fig:convergence}B). Among its $74$ members, the strongest
LDL-C contributions come from the receptor-uptake arm---\textit{LDLR}
($p=3.1\times10^{-24}$), \textit{PCSK9} ($7.8\times10^{-23}$),
\textit{LDLRAP1} ($7.4\times10^{-22}$), with \textit{APOB}, \textit{APOE} and
\textit{MTTP} close behind---while for HDL-C they come from the reverse-transport
and lipolysis arm: \textit{ABCA1} ($1.3\times10^{-33}$), \textit{LIPC}
($5.8\times10^{-29}$), \textit{SCARB1} ($1.9\times10^{-22}$), \textit{LPL}
($8.9\times10^{-22}$), with \textit{APOA1}, \textit{ANGPTL3} and \textit{LMF1}
following. Only five of the fifteen leading drivers are common to both traits.

This is the behaviour a pathway method exists to produce. Two GWAS that
disagree about which genes matter are nonetheless aggregated onto the same
biological process, because that process contains both arms. The convergence is
not a restatement of a shared gene-level signal; the gene-level signals are
largely disjoint.

Two qualifications belong with the claim. First, the two traits are not
independent replicates at the pathway level: across all $1{,}390$ sets the
Spearman correlation of the pathway statistic was $0.863$, and $75$ of the top
hundred pathways were shared (Figure~\ref{fig:convergence}C). That correlation
is driven by the size-dependent background of Section~3.9.1, which is common to
both traits, rather than by the leading hits---which is precisely why the
driver-level comparison, not the aggregate correlation, is the evidence for the
claim above. Second, the phenotypes themselves are genetically correlated to a
degree we have not estimated here; the gene-level Spearman of $0.393$ is the
quantity that actually enters CATFISH and is reported in its place.

Below the first rank, large generic sets---sumoylation, Golgi traffic, cellular
response to chemical stress, rRNA processing---are interleaved with the lipid
pathways in both traits, again consistent with the size gradient. Lipid,
lipoprotein and sterol sets nonetheless concentrate near the top of both
orderings, and they move between traits in a biologically sensible way:
PPAR$\alpha$-regulated lipid metabolism ranks $16$th for LDL-C and $7$th for
HDL-C, while sphingolipid metabolism ranks $81$st and $182$nd respectively.

\begin{figure*}[!t]
\centering
\maybefig{Fig10_ldl_hdl_convergence.png}{\textwidth}{LDL-C and HDL-C reach the
same pathway from different genes. \textbf{A}~Gene-level $-\log_{10}p$ for the
$18{,}563$ genes analysed in both traits; the $74$ members of the top-ranked
pathway are highlighted and selected drivers labelled. The two GWAS agree only
moderately (Spearman $0.393$; $13$ of the top hundred genes shared).
\textbf{B}~Within that pathway, the twelve strongest contributors in each
trait. LDL-C is carried by the receptor-uptake arm and HDL-C by the
reverse-transport and lipolysis arm, with five of fifteen leading drivers in
common. \textbf{C}~Pathway ranks in the two traits, lipid-related sets
highlighted and the shared top pathway circled; it ranks first in both. The
broad agreement away from the top reflects the size-dependent background of
Section~3.9.1 rather than shared leading signal.}{fig:convergence}
\end{figure*}

\subsubsection*{3.9.3 The adaptive TFisher component is defensible but not
load-bearing}

The analyses above use all five component tests. Because the soft TFisher
component is adaptive over a grid of truncation thresholds and is by some
margin the slowest to evaluate, we also ran the complete HDL-C analysis with it
removed, at matched $B=10{,}000$, as a direct ablation.

The two orderings are almost identical: the Spearman rank correlation between
the four- and five-component results was $0.9994$, the median rank change was
three places, and the largest change over all $1{,}390$ pathways was $104$
places, with $44$ sets moving more than fifty. Every lipid pathway moved
slightly upward when TFisher was included---plasma lipoprotein assembly
$2\rightarrow1$, PPAR$\alpha$ $8\rightarrow7$, lipoprotein clearance
$98\rightarrow87$. This matches the leave-one-out simulation of Section~3.4,
where removing TFisher left the family-wise error rate unchanged at $0.042$.

We report this because the adaptive grid is the component most open to the
objection that taking a minimum across five thresholds without correcting for
the search is anti-conservative. That objection is correct for the analytic
component---its uncalibrated type~I error reaches $\nCmpTFisherStr$ at strong
LD---but the resampling applies the same grid search to every null draw, so the
inflation cancels, and the calibrated component is the best behaved of the five
($\nCmpTFisherCal$). The ablation confirms that nothing in the real-data
conclusions depends on it either way.

\subsubsection*{3.9.4 Gene prioritization is unaffected}

The three-layer score of Section~2.5 was computed over all $\nCandN$ genes.
Table~\ref{tab:candidates} gives the top five and Supplementary
Table~\ref{tab:candidates_supp} the top fifty.

\textit{LDLR} ranked $\nCandLDLR$ and \textit{PCSK9} $\nCandPCSKnine$---the
receptor that clears LDL from plasma and the protease that degrades it, the two
genes around which LDL-lowering pharmacology is built. \textit{APOE} and
\textit{APOC1}, at the chromosome~19 apolipoprotein cluster, follow close
behind. Three of the top five are therefore genes with direct causal roles in
LDL-C biology.

The other two require care, because they illustrate what a gene-based test can
and cannot do. \textit{PVRL2} (chr19:45.35--45.39\,Mb) lies $17\,\mathrm{kb}$
from \textit{APOE}, and \textit{SMARCA4} (chr19:11.07--11.17\,Mb) lies
$27\,\mathrm{kb}$ from \textit{LDLR}; \textit{KIAA1324}
(chr1:109.66--109.75\,Mb), just outside the top five, sits in the
\textit{SORT1}/\textit{CELSR2}/\textit{PSRC1} region. None of these is a
lipid gene. Each is an LD neighbour of one, and ranks highly because the
association signal at a locus is shared across the genes within it. This is the
expected behaviour of every gene-based test, CATFISH included: the method ranks
\emph{loci}, and resolving the causal gene inside a locus requires evidence it
does not use. We state this explicitly because a table of top-ranked genes
invites the opposite reading.

Figure~\ref{fig:candlayers} shows where each of the top twenty genes sits in
each layer, which separates genes carried by a single layer from those with
concordant support across all three.

As a pre-specified check, all thirteen established LDL-C genes named in
Section~2.7 were recovered, with a median rank of $\nCandCtrlMedian$ out of
$\nCandN$ against a random expectation of $\nCandRandom$, and
$\nCandCtrlTopPct$ of $\nCandCtrlTotal$ in the top $1\%$.

The informative detail is \emph{which} layer produces this. The pathway ranks of
the control genes are good but unremarkable---\textit{LDLR} ranks $462$nd of
$\nCandNPath$ pathway-annotated genes---while their GWAS and MAGMA ranks are
near the top. The ordering is driven by the gene-level layers, and the pathway
layer contributes little. That is the saturation result of Section~3.9.1 seen
from the other side: saturation degrades the pathway layer specifically, because
the self-contained hypothesis has stopped discriminating between gene sets,
while leaving gene-level resolution intact. A composite score that weights
pathway evidence at half the strength of direct genetic evidence degrades
gracefully in exactly this regime, which was the reason for the weighting
(Section~2.5.3) but not a situation it was designed against.

One implementation point affects this table directly and generalises beyond it.
When the GWAS layer is ranked on probabilities reconstructed from GLGC and
floored at $10^{-300}$, $\nCandGwasTies$ genes tie at the floor---including
\textit{LDLR}, \textit{PCSK9}, \textit{APOB}, \textit{APOE}, \textit{LPA} and
\textit{HMGCR}---and the layer cannot order the very genes it most needs to.
Ranking on unfloored $-\log_{10}p$ removes the tie at no other cost. This is the
same double-precision limit as Section~2.4, surfacing in a different part of the
pipeline, and it is worth checking wherever a ranking is built from
large-$n$ GWAS summary statistics.

\maybeinput{auto/tab_candidates.tex}

\begin{figure}[!t]
\centering
\maybefig{Fig8_candidate_layers.png}{\columnwidth}{Evidence layers for the
twenty highest-scoring genes in the LDL-C analysis. Each row is a gene; the
three points give its rank within the locus-level GWAS layer, the MAGMA
gene-level layer and the CATFISH pathway layer, on a log rank scale with one
the strongest. Genes in bold are the pre-specified positive controls. The
genetic layers place the controls near the top while their pathway ranks are
unremarkable, which is the saturation of Section~3.9.1 seen gene by
gene.}{fig:candlayers}
\end{figure}

\begin{table}[!t]
\centering
\caption{Gene-level signal density and pathway-size dependence across the four
GWAS analysed. The last column is the Spearman correlation between pathway size
and the pathway-level evidence ($-\log_{10}$ omnibus $p$); more negative values
indicate that larger pathways receive systematically stronger evidence.
n.c., not computed.}
\label{tab:density}
\small
\begin{tabular}{lrrr}
\hline
\textbf{Dataset} & \makecell[r]{\textbf{Genes GW}\\\textbf{significant}} &
\makecell[r]{\textbf{Median}\\\textbf{gene $p$}} &
\makecell[r]{\textbf{Size}\\\textbf{correlation}} \\
\hline
Maize grain N (MLM)        & 0.00\%  & 0.468 & $-0.179$ \\
Sorghum stem volume        & 0.86\%  & 0.440 & n.c.     \\
Maize grain N (GLM)        & 1.15\%  & 0.185 & $-0.667$ \\
Human LDL-C (GLGC, EUR)    & 11.2\%  & 0.083 & $-0.807$ \\
\hline
\end{tabular}
\end{table}

\subsubsection*{3.9.5 Signal density across datasets}

Table~\ref{tab:density} places the two real analyses in context. Gene-level
signal density varies by more than an order of magnitude between the maize MLM
analysis and the human meta-analysis, and the dependence of the pathway
statistic on pathway size tracks it closely, from $-0.179$ in the sparsest
dataset to $-0.807$ in the densest. This is not a bias to be corrected away. In
a saturated regime larger sets genuinely do accumulate more association signal,
because the background itself is enriched; the statistic reports something true
about the data and uninteresting about the biology.

We note in passing that the maize GLM analysis, included here for contrast,
shows a median gene-level $p$-value of $0.185$, far below the $0.5$ expected
under the null. This is population-structure inflation rather than signal, and
is why the MLM analysis is the one we report in Section~3.8. It also
illustrates how readily a saturated-looking pathway result can arise from a
mis-specified gene-level model rather than from biology, which is an argument
for inspecting the gene-level null before interpreting any pathway output.

%%%%%%%%%%%%%%%%%%%%%%%%%%%%%%%%%%%%%%%%%%%%%%%%%%%%%%%%%%%%%%%%%%%%%%%
\section{4 Discussion}

\subsection*{4.1 What the Westfall--Young layer buys, and what it does not}

The contribution of the Westfall--Young layer is narrow and should be stated
narrowly. Reading the null matrix across its rows rather than down its columns
costs nothing---the draws have already been taken---and it makes the family-wise
statement exact under the actual overlap structure of the pathway database being
used, without requiring that structure to be estimated or even described. It
does not change any pathway's marginal evidence, which is what the column-wise
reading gives, and it does not improve power. In the shared-universe simulations with no within-pathway correlation, where
the analytic $p$-value is valid, the Westfall--Young threshold was $\nWyGainLo$
to $\nWyGainHi$ times the Bonferroni threshold $\nBonfThrK$: modestly more
permissive, which is the gain from accounting for overlap at an average of
$\nBrkMembership$ pathways per gene. That is enough to matter for a borderline result and not enough to
rescue an analysis that has gone wrong for other reasons.

The practical value is diagnostic as much as inferential. The ratio of the
Bonferroni threshold to the realized Westfall--Young threshold is a single
number that tells an analyst how far from independent their pathway family
actually is. A ratio near one says the database behaves like a set of
independent tests and Bonferroni is close to right; a ratio far from one says
the correction in common use is substantially miscalibrated for this particular
collection of gene sets. That number is cheap to compute as a by-product of the
procedure and, as far as we are aware, is not routinely reported.

The ratio has to be read on calibrated statistics to mean this. In the same
simulations with within-pathway correlation, the threshold set on the analytic
scale was $\nWyStrictLo$ to $\nWyStrictHi$ times \emph{stricter} than
Bonferroni---not because the pathway family was more dependent, but because the
analytic $p$-value was anti-conservative and the Westfall--Young threshold,
drawn from the same null, corrected for it. On that scale the ratio mixes the
miscalibration of the per-test statistic with the dependence among tests and
cannot separate them.

We were initially wrong about one aspect of this and correct the record here. We
expected the size-standardized Westfall--Young variant of Section~2.3 to remove
the dependence of pathway evidence on pathway size, on the reasoning that
expressing each statistic against its own null bank puts all pathways on a
common scale. It did not: in the maize data the size--evidence correlation moved
from $-0.807$ to $-0.800$ in the human analysis and remained essentially
unchanged elsewhere. The reason is visible once measured. The standard
deviations of the per-size null banks differ by only a few percent across the
pathway sizes we examined, so the standardization divides by a nearly constant
factor and cannot remove a dependence that lives in the numerator. The size
dependence is in the observed statistics, and in a saturated regime it is real
accumulation of signal, not a scaling artefact. Standardization is still worth
reporting because it changes individual rankings, but it is not a correction for
size bias and should not be presented as one.

\subsection*{4.2 Saturation and the self-contained hypothesis}

The LDL-cholesterol result in Section~3.9 is, we think, the most consequential
finding in this paper, and it is not about CATFISH. Self-contained gene-set
tests ask whether the genes in a set show more association than the null
distribution of no association \citep{Goeman2007,Maciejewski2014}. That question
has a definite answer, and at $n\approx1.3$ million for a highly polygenic trait
the answer is yes for almost every set of genes, because $11.2\%$ of all genes
are individually genome-wide significant. A random set of fifty genes rejects
the self-contained null with probability $0.997$. Nothing has malfunctioned; the
hypothesis being tested has simply stopped discriminating.

Three consequences follow. First, as GWAS sample sizes grow, self-contained
testing will fail on more traits, and it will fail silently, producing long
lists of highly significant pathways that look like results. Second, no amount
of multiple-testing sophistication fixes this: we applied exact family-wise
control and $1{,}374$ of $1{,}390$ pathways remained significant. Reporting a
more careful threshold on a saturated analysis produces a more carefully wrong
answer. Third, the competitive hypothesis---whether the genes in a set are
\emph{more} associated than genes outside it---does not saturate in the same
way, because the comparison is against the observed genome-wide background
rather than against no association. For large, well-powered human GWAS, the
competitive framing is the appropriate one, and the self-contained framing that
CATFISH implements is best suited to the low-signal, strongly polygenic regime
illustrated by the maize analysis, where no gene reaches genome-wide
significance and a ranking is the only available output.

This is a boundary on the applicability of the method we are presenting, and we
would rather state it than discover it in review. It also explains an initially
puzzling observation: the pathways CATFISH ranked highest in the human analysis
were things like spliceosome assembly and the TCA cycle, which have no
particular relationship to lipid metabolism but are large and contain many of
the genes that are significant in a saturated genome.

\subsection*{4.3 Numerical precision is a correctness issue}

Section~2.4 reads like an implementation note and is not one. Had we not moved
to log space, this paper would have reported a top-pathway list that was
determined by file order, with every leading $p$-value equal to the
double-precision floor, and nothing in the output would have indicated that
anything was wrong. We encountered exactly this failure twice during
development, once from ACAT overflow and once from underflow of the combined
$p$-value, and in both cases the symptom was a plausible-looking table.

The simulations produced a third instance, in a different component and with a
different symptom. The soft TFisher implementation returns the null
distribution function of its statistic, and the $p$-value is one minus that
value; past $1-10^{-16}$ the difference is exactly zero. In the exchangeable
design at $\rho=0.6$, $m=50$, about $3.5\%$ of null pathways reached that point
and were reported at the general lower floor. With the floor at $10^{-50}$ the
empirical family-wise error rate in that configuration was $0.24$. Lowering the
floor to $10^{-300}$, in an attempt to recover resolution, raised it to $1.00$:
the spurious value became more extreme than anything another component
produced, so the omnibus family minimum was the same shared number in
essentially every simulated family and the Westfall--Young threshold sat on the
tie. The fix was to cap the component at the precision it has. At $10^{-16}$,
the components with closed-form tails---which in the same configuration resolve
genuinely distinct values below $10^{-100}$---determine the family minimum, and
it is once again a continuous statistic. A floor is a claim about resolution; a
floor set below the resolution a computation actually has manufactures ties that
look like extreme evidence.

We therefore suggest that any pathway method intended for use on modern
large-$n$ GWAS summary statistics should state explicitly how it represents
$p$-values below $10^{-300}$, and that reviewers should ask. The arithmetic is
not difficult; the problem is that the failure is invisible unless specifically
checked for. A simple diagnostic is to count how many pathways share the
smallest reported $p$-value. If that count exceeds one or two, the ranking below
that point is not meaningful. For a family-wise procedure the corresponding
check is the number of distinct values taken by the null family minimum; a
threshold drawn from a distribution with a large atom is not a threshold.

\subsection*{4.4 Practical recommendations}

For users, our recommendations are as follows. Inspect the gene-level null
before running any pathway analysis: a median gene-level $p$-value far from
$0.5$ indicates either genuine saturation or population-structure inflation, and
these require different responses (the former argues for competitive testing,
the latter for a mixed model). Calibrate against the correlated null whenever
gene--gene correlations are available, and verify that the correlation file
covers all chromosomes; in one of our own runs only chromosome~1
correlations had been parsed, and $85\%$ of the genome was silently treated as
independent, a failure that produced no error and no obviously wrong output.
Report the per-pathway calibrated $p$-values and the Westfall--Young family-wise
ones, together with the Bonferroni threshold and its ratio to the
Westfall--Young threshold computed on calibrated statistics (Section~4.1), since
the comparison is informative about the database. Before trusting any
family-wise threshold, count the distinct values taken by the null family
minimum. When using adaptive soft TFisher, choose a truncation grid that spans
small and large thresholds; a grid confined to small $\tau$ leaves the statistic
empty on most small pathways. Finally, in the low-signal regime, treat the ranking rather than the
significance list as the deliverable, and validate it against independent
biology where possible, as we have attempted to do with the maize analysis.

%%%%%%%%%%%%%%%%%%%%%%%%%%%%%%%%%%%%%%%%%%%%%%%%%%%%%%%%%%%%%%%%%%%%%%%
\section{5 Limitations and future work}

Several limitations bound what can be concluded from this work, and we list them
explicitly.

\textbf{No simulated GWAS.} The simulations in Section~2.6 operate at the level
of the omnibus statistic and its null, using block-structured and exchangeable
correlation models for gene--gene dependence. We did not simulate genotypes, assign causal
variants, generate phenotypes, and push them through the full
GWAS--MAGMA--CATFISH pipeline. That experiment is the natural next step and is
the one that would establish power against realistic genetic architectures
rather than against latent mean shifts; it was deliberately excluded here rather
than attempted partially. Consequently, the archetype power curves of Section~3.7
are power against stylised signal shapes injected on the gene-level $z$ scale,
not against realistic genetic architectures.

\textbf{Unreported correlations are set to zero.} The correlation matrices of
Section~2.3 are built by filling MAGMA's reported gene pairs into an identity
matrix, so any pair MAGMA does not report is treated as independent. For distant
gene pairs this is close to true and the matrix is banded by construction, but a
nearby pair that goes unreported---because it fell outside a batch boundary, or
because one gene was dropped---becomes a genuine understatement of dependence,
in the anti-conservative direction. Section~3.6 quantifies the cost, and
corrects an assumption we made before measuring it: the two readings do
\emph{not} inherit the limitation equally. The family-wise reading degrades
faster than the marginal one, because the family minimum is a tail
statistic and an error in the correlation structure does its worst damage in
the tail. Verifying chromosome coverage of the correlation file catches the
gross form of this failure; \texttt{10\_pair\_coverage.R} measures the
fine-grained form, and we recommend reporting it alongside any family-wise
$p$-value.

\textbf{Simplified LD models.} Real gene--gene correlation is driven by physical
proximity and decays continuously with distance. The shared-universe simulations
approximate this with independent blocks of exchangeable correlation, so a
pathway spanning several blocks has mostly uncorrelated members; the
exchangeable design instead correlates every member of a pathway, which is far
stronger than block-structured LD produces and is used only as a stress test.
Neither reproduces the gradual decay of correlation with distance. The real-data
analyses use empirical MAGMA correlations and are not subject to this
limitation.

\textbf{What strong control rests on.} The argument of Section~2.3 for strong
family-wise control is exact under its model and inherits that model's limits.
It assumes association acts as a shift in the means of associated genes'
$z$-scores, leaving the correlation among unassociated genes unchanged, and it
defines a null pathway as one containing no associated gene. Under linkage
disequilibrium a gene can be associated at the gene level solely because it is
in LD with a causal variant in a neighbouring gene; a pathway containing it is
then non-null under the self-contained gene-level hypothesis even if it contains
no causal gene, so rejecting it is not a family-wise error, although it may be a
false positive biologically. The argument also inherits the adequacy of the
correlation matrix, and Section~3.6 shows how quickly family-wise calibration
degrades when that matrix is incomplete. We have not separately confirmed strong
control by simulation; under the model it follows from the complete-null
results.

\textbf{Parametric null.} The null draws come from a multivariate normal model
of gene--gene correlation rather than from permuting phenotypes. The procedure
is Westfall--Young in construction, but its validity is that of the correlation
model, which Section~3.6 makes concrete.

\textbf{Soft TFisher resolution.} With the current implementation the soft
TFisher $p$-value cannot be resolved below about $10^{-16}$, so among strongly
enriched pathways that component ties. The other components resolve far beyond
this and determine the omnibus ordering there; a log-scale implementation of
the soft-thresholding tail would remove the limitation.

\textbf{No competitive arm.} We were unable to include a direct comparison
against MAGMA's competitive gene-set test in this version, because the
competitive entry point is not currently exported from the CATFISH package
namespace. Given the argument of Section~4.2, that comparison is the single most
important addition to make, and it is in progress.

\textbf{Incomplete cross-species replication.} The Westfall--Young analysis was
run on maize and human data. The sorghum stem-volume dataset used in the
original CATFISH paper could not be re-analysed under the new calibration
because the MAGMA \texttt{.genes.raw} file, which the genome-wide correlated
draw requires, was not available in the archived results; its signal-density
summary in Table~\ref{tab:density} is reported from the gene-level output alone.

\textbf{Tail resolution.} Monte Carlo calibration resolves $p$-values only to
$1/(B+1)$. CATFISH offers a generalized Pareto tail extrapolation for the
extreme tail \citep{Knijnenburg2009}, but it requires a minimum number of valid
null draws, and a single missing draw can silently disable it; in one of our
human runs this caused $638$ of $1{,}348$ pathways to tie at the Monte Carlo
floor. The threshold is now set to half the number of permutations rather than
to the full number, but we have not yet verified the fix at production scale.

\textbf{Dependence on upstream choices.} Results remain sensitive to the
SNP-to-gene window, the reference panel used for LD, gene model definitions, and
the pathway database, as is true of every method in this class. The $10\,$kb
windows used here are conventional but not optimal for any particular trait.

%%%%%%%%%%%%%%%%%%%%%%%%%%%%%%%%%%%%%%%%%%%%%%%%%%%%%%%%%%%%%%%%%%%%%%%
\section{6 Conclusion}

Pathway analysis of GWAS summary statistics faces two distinct dependence
problems, and the literature has concentrated on the easier one. Correlation
among genes within a pathway is widely acknowledged and routinely modelled;
correlation among pathways, induced by the shared gene membership that is
intrinsic to how pathway databases are built, is usually handled with a
correction that assumes it away. We have added a Westfall--Young max-statistic
layer to CATFISH that controls the family-wise error rate in the strong sense
whatever that dependence, under the same correlation model used for
calibration, requires nothing to be estimated about it beyond
the gene--gene correlation already used for calibration, and yields as a
by-product a direct measure of how far Bonferroni is from correct for any given
database, provided it is read on calibrated statistics.

Two further findings matter beyond this particular method. Computing the omnibus
in log space is a correctness requirement at modern GWAS sample sizes, not an
optimisation, and the failure it prevents is silent. And self-contained pathway
testing saturates: in a GWAS of $1.3$ million individuals, $11.2\%$ of genes are
genome-wide significant, a random fifty-gene set rejects the self-contained null
with probability $0.997$, and exact family-wise error control leaves $1{,}374$
of $1{,}390$ pathways significant. That is a property of the hypothesis, not of
the correction. CATFISH, and self-contained pathway testing generally, is best
suited to the low-signal polygenic regime, where the maize grain-nitrogen
analysis recovered L-homoserine biosynthesis as the leading pathway and a known
aspartate kinase isozyme as the leading gene out of $34{,}873$, in agreement
with independent prior work.

%%%%%%%%%%%%%%%%%%%%%%%%%%%%%%%%%%%%%%%%%%%%%%%%%%%%%%%%%%%%%%%%%%%%%%%
\section{7 Conflicts of interest}
The authors declare that they have no competing interests.

\section{8 Funding}
This research was supported by the DOE Office of Science, Office of Biological
and Environmental Research (BER), grant no. DE-SC0021889 for North Carolina
State University. This program is supported by the U.S. Department of Energy,
Office of Science, through the Genomic Science Program, Office of Biological and
Environmental Research, under FWP ERKP123.

\section{9 Data availability}
The GLGC LDL-cholesterol summary statistics are available from the Global Lipids
Genetics Consortium. The CATFISH R package and all code used to generate the
figures, tables and simulations underlying this article are available at
\href{https://github.com/nirwantandukar/CATFISH}{github.com/nirwantandukar/CATFISH},
in the \texttt{catfish2/} directory.

\section{10 Author contributions statement}
N.T. conceived and implemented the method, designed and ran the simulations,
analysed the real data, and wrote the manuscript. J.-Y.T. advised on the
statistical framework and calibration. R.R.-A. supervised the work and revised
the manuscript. All authors read and approved the final manuscript.

\section{11 Acknowledgments}
The authors thank the anonymous reviewers for their valuable suggestions.

\clearpage
\section*{Supplementary material}

\maybeinput{auto/tab_candidates_supp.tex}

\clearpage
\bibliographystyle{oup-abbrvnat}
\bibliography{refs}

\end{document}
